\ECHead{Proofs}
\section{Coordinated Control Algorithm}
\label{app:algorithm}

Algorithm~\ref{alg:coordinated_control} details the coordinated threshold adjustment procedure described in Section~\ref{subsec:coordinated_control}.
At each iteration, the algorithm estimates the steady-state system variables $(z_0^*, z_{1i}^*, z_{2i}^*, q_i^*)$ by time-averaging observations from the matching system, computes the current KMI values $\boldsymbol{\zeta}^{(k)}$ using the closed-form expressions~\eqref{eq:kmi_welfare} or~\eqref{eq:kmi_revenue}, and estimates the Jacobian via finite differences.
It then inverts the $2 \times 2$ Jacobian matrix and applies the Newton update with projection onto $\boldsymbol{\mu} \geq \mathbf{0}$.

\begin{algorithm}
\caption{Coordinated Threshold Adjustment (Newton--Raphson Type)}
\label{alg:coordinated_control}
\begin{algorithmic}[1]
\State \textbf{Initialize:} $\boldsymbol{\mu}^{(0)}$, tolerance $\epsilon$, step size $\eta$.
\While{$\|\boldsymbol{\zeta}(\boldsymbol{\mu}^{(k)}) - \mathbf{1}\| > \epsilon$}
    \State Estimate steady-state variables $(z_0,z_{1i},q_i)$ by time averaging.
    \State Compute $\boldsymbol{\zeta}^{(k)}$ via \eqref{eq:kmi_welfare} or \eqref{eq:kmi_revenue}.
    \State Estimate Jacobian $\hat{\mathbf{J}}$ via finite differences:
    \begin{equation*}
      \hat{J}_{ij}
      \approx
      \frac{\zeta_i(\mu_{1j}+\delta)-\zeta_i(\mu_{1j})}{\delta}.
    \end{equation*}
    \State Compute direction
    \begin{equation*}
      \mathbf{d}_k
      =
      -\hat{\mathbf{J}}^{-1}
      (\boldsymbol{\zeta}^{(k)}-\mathbf{1}).
    \end{equation*}
    \State Update
    $\boldsymbol{\mu}^{(k+1)}
    \leftarrow
    \max(\mathbf{0},\,\boldsymbol{\mu}^{(k)}+\eta\mathbf{d}_k)$.
    \State $k\leftarrow k+1$.
\EndWhile
\end{algorithmic}
\end{algorithm}

The computational burden is minimal: each Jacobian estimation requires only two additional evaluations of the two-dimensional fluid system (one per threshold direction), and the matrix inversion involves only a $2 \times 2$ system.
The entire update step executes in sub-second time, making the algorithm suitable for online deployment.
The algorithm handles both welfare and revenue objectives by selecting the appropriate KMI formula.

\section{Experiment Implementation Details and Additional Results}
\label{app:experiment_details}

\subsection{Baseline Parameters}
Unless otherwise noted, the numerical experiments utilize the baseline parameters detailed in Table~\ref{tab:baseline_params}. The system operates in a stylized two-tier market with driver supply $K=1000$.

\begin{table}[htbp]
\centering
\caption{Baseline Simulation Parameters}
\label{tab:baseline_params}
\small
\begin{tabular}{lll}
\toprule
Parameter & Symbol & Value \\
\midrule
Driver Supply & $K$ & 1000 \\
Arrival Rates (Balanced) & $\lambda_1, \lambda_2$ & $2.0K, 3.0K$ \\
Prices & $p_1, p_2$ & 2, 1 \\
Pre-match Patience & $\theta_{01}, \theta_{02}$ & 2, 1 \\
Post-match Patience & $\theta_{11}, \theta_{12}$ & 3, 4 \\
Spatial Parameters & $C$ & 50 \\
Matching Exponents & $\alpha_1, \alpha_2$ & 0.5, 0.5 \\
Trip Completion Rate & $\mu_2$ & 10 \\
\bottomrule
\end{tabular}
\end{table}

\subsection{3D Objective and KMI Landscapes}
\label{app:3d_landscapes}
To complement the contour plots in Section~\ref{sec:validate}, we provide 3D visualizations of both the objective functions and the KMI surfaces.

Figure~\ref{fig:3d_objectives} illustrates the anisotropy of the optimization landscape, highlighting why direct gradient ascent on thresholds can be unreliable.
Figure~\ref{fig:3d_kmi} displays the KMI surfaces, demonstrating how the $\zeta$ indices provide a regularized gradient field that robustly guides the system toward the optimum.

\begin{figure}[htbp]
\centering
\caption{3D Visualization of Welfare and Revenue Landscapes}
\label{fig:3d_objectives}
\begin{subfigure}{0.48\linewidth}
\centering
\captionsetup{font=normalsize}
    \includegraphics[width=\linewidth]{outputs/exp2_W/exp2_thresholds_busy_3D.pdf}
    \caption{Welfare Surface}
\end{subfigure}\hfill
\begin{subfigure}{0.48\linewidth}
\centering
\captionsetup{font=normalsize}
    \includegraphics[width=\linewidth]{outputs/exp2_R/exp2_thresholds_busy_3D.pdf}
    \caption{Revenue Surface}
\end{subfigure}
\end{figure}

\begin{figure}[htbp]
\centering
\caption{3D Visualization of Key Matching Indices (KMI)}
\label{fig:3d_kmi}
\begin{subfigure}{0.48\linewidth}
\centering
\captionsetup{font=normalsize}
    \includegraphics[width=\linewidth]{outputs/exp2_W/exp2_kmi_busy_3D.pdf}
    \caption{Welfare KMI Surface ($\zeta^W$)}
\end{subfigure}\hfill
\begin{subfigure}{0.48\linewidth}
\centering
\captionsetup{font=normalsize}
    \includegraphics[width=\linewidth]{outputs/exp2_R/exp2_kmi_busy_3D.png}
    \caption{Revenue KMI Surface ($\zeta^R$)}
\end{subfigure}
\end{figure}

\subsection{Convergence Rate and Hyperparameter Tuning}
\label{app:tuning}

%We analyze the convergence rate as a function of the learning rate $\eta$. 
%The performance metric is the \emph{time-to-$\varepsilon$}, defined as the 
%time required for the combined log-error $\|\log \boldsymbol{\zeta}\|_2$ to 
%drop below $\varepsilon=0.01$.
%
%Figure~\ref{fig:exp3-eta} reports the mean time-to-$\varepsilon$ (averaged 
%over 100 replications). The data reveal a convex trade-off. Small $\eta$ results 
%in sluggish convergence due to insufficient step sizes. Conversely, excessively 
%large $\eta$ degrades performance, as evidenced by widening confidence intervals 
%that signal oscillatory behavior. An \textbf{effective operating region} 
%(e.g., $\eta \approx 0.15$) exists where the controller minimizes convergence 
%time while maintaining low variance.
\subsubsection{Convergence Rate Analysis}
To guide practical implementation, we analyze the controller's convergence rate as a function of the learning rate, \(\eta\).
We define the performance metric as the \textbf{time-to-\(\varepsilon\)}, the time required for the combined log-error, \(\sqrt{(\log\zeta_1)^2 + (\log\zeta_2)^2}\), to fall below a threshold of \(\varepsilon=0.01\).
For statistical robustness, we average results over 100 replications for each \(\eta\) value.

Figure~\ref{fig:exp3-eta} plots the mean time-to-\(\varepsilon\) against \(\eta\), with error bars representing 95\% confidence intervals.
The results reveal a distinct trade-off that defines an optimal tuning window.
For small \(\eta\), convergence is slow, as the small step size leads to insufficient update magnitudes.
As \(\eta\) increases, the convergence time drops sharply, identifying a "sweet spot" where the controller is both fast and stable, evidenced by narrow confidence intervals that indicate low variance across replications.
Beyond this region, further increases in \(\eta\) yield diminishing returns; the curve flattens and eventually rises, while the widening confidence intervals signal growing instability and oscillatory behavior near the stability boundary.
This analysis provides an actionable guide for selecting an effective \(\eta\) that balances convergence speed with robustness.

\begin{figure}[htbp]
\centering
\caption{Convergence rate diagnostics}
\label{fig:exp3-eta}
\includegraphics[width=0.5\linewidth]{outputs/exp3/exp3_eta_vs_time_to_eps.pdf}
\begin{minipage}{0.9\textwidth}
    \footnotesize\raggedright
    \textit{Notes.} Mean time-to-$\varepsilon$ vs. learning rate $\eta$ (with 95\% confidence intervals).
\end{minipage}
\end{figure}

%Figure~\ref{fig:high-eta} further illustrates the stability boundary. 
%An aggressive rate ($\eta=0.30$) induces overshoot and sustained oscillations 
%in the log-error, whereas a conservative rate ($\eta=0.04$) ensures smooth but 
%slower convergence.
\subsubsection{Stability Boundary and Implementation}
The choice of learning rate \(\eta\) is critical for practical deployment.
While a larger \(\eta\) can hasten convergence, an excessively high value may lead to instability.
Figure~\ref{fig:high-eta} demonstrates this by comparing a conservative rate (\(\eta=0.04\)) with an aggressive one (\(\eta=0.30\)).
The aggressive rate causes the thresholds to overshoot the optimal values, and the log-error exhibits oscillations, indicating proximity to the stability boundary.
This suggests that \(\eta\) should be tuned to balance speed and stability, possibly with an adaptive schedule that reduces the rate as the system approaches the optimum.

From an implementation perspective, the \(\zeta\)-guided controller is efficient.
It relies on summary statistics that are typically available in real-time platform monitoring systems.
The Jacobian of the \(\zeta\) mapping can be estimated infrequently using small perturbations, while the thresholds are updated more frequently.
This two-timescale approach ensures that the computational overhead remains low, making the controller suitable for large-scale, dynamic environments.

\begin{figure}[htbp]
\centering
\caption{Stability boundary demonstration}
\label{fig:high-eta}
\begin{subfigure}{0.48\linewidth}
\centering
\captionsetup{font=normalsize}
    \includegraphics[width=\linewidth]{outputs/exp3/exp3_high_eta_thresholds.pdf}
    \caption{Threshold Overshoot}
\end{subfigure}\hfill
\begin{subfigure}{0.48\linewidth}
\centering
\captionsetup{font=normalsize}
    \includegraphics[width=\linewidth]{outputs/exp3/exp3_high_eta_log_error.pdf}
    \caption{Log-Error Oscillations}
\end{subfigure}
\end{figure}

\section{Proof of Results}

\subsection{Proof of Theorem \ref{thm:general_optimality} and Proposition \ref{prop:first_level}}
\label{sec:ec_optimality}

This section derives the first-order necessary conditions for the optimization problem \eqref{eq:opt_problem}.
The objective is to maximize the weighted throughput $J(\boldsymbol{\mu}) = w_1 z_{21}^* + w_2 z_{22}^*$ subject to the fluid equilibrium constraints derived in Section 3.

\subsubsection{Lagrangian Formulation}
Let $\lambda_i, \theta_{0i}, \theta_{1i}, \mu_{2}$ be the system parameters.
The equilibrium is governed by the flow balance and state constraints.
We construct the Lagrangian $\mathcal{L}$ associated with the optimization problem.
Note that while the fluid dynamics involve $z_{1i}$ and $q_i$, the control variables are the thresholds $\mu_{1i}$ which impose the constraint $C(q_i^*)^{\alpha_1}(z_0^*)^{\alpha_2} = \mu_{1i}^*$.
For simplicity, we use the notation $(z_{2i}, z_{1i}, q_i, z_0; i=1,2)$ to denote the equilibrium state variables $(z_{2i}^*, z_{1i}^*, q_i^*, z_0^*; i=1,2)$ in the remainder of Section~\ref{sec:ec_optimality}.

The Lagrangian is
\begin{align}
\mathcal{L} &= w_1 z_{21} + w_2 z_{22} \nonumber \\
&\quad + \sum_{i=1}^{2} \beta_i (\lambda_i - q_i \theta_{0i} - z_{1i} \theta_{1i} - z_{2i} \mu_2) \quad \text{(Passenger Flow Balance)} \nonumber \\
&\quad + \sum_{i=1}^{2} \gamma_i \left( z_{1i}\mu_{1i} - z_{2i}\mu_2 \right) \quad \text{(Service Rate Balance)} \nonumber \\
&\quad + \gamma_3 \left(1 - z_0 - \sum_{i=1}^2 (z_{1i} + z_{2i}) \right) \quad \text{(Driver Conservation)}. \label{eq:Lagrangian}
\end{align}
Here, $\beta_i, \gamma_i, \gamma_3$ are the Lagrange multipliers (shadow prices) associated with the respective constraints.

\subsubsection{First-Order Conditions (FOCs)}
Assuming an interior solution exists where all state variables are strictly positive, the Karush-Kuhn-Tucker (KKT) conditions require the partial derivatives of $\mathcal{L}$ with respect to the state variables $(z_{2i}, z_{1i}, q_i, z_0)$ to satisfy stationarity:

\begin{align}
\frac{\partial \mathcal{L}}{\partial z_{2i}} &= w_i - \beta_i \mu_2 - \gamma_i \mu_2 - \gamma_3 = 0 \implies \gamma_3 = w_i - (\beta_i + \gamma_i)\mu_2, \quad i \in \{1,2\} \label{eq:foc_z2} \\
\frac{\partial \mathcal{L}}{\partial z_{1i}} &= -\beta_i \theta_{1i} + \gamma_i\mu_{1i} - \gamma_3 = 0 \implies \gamma_i\mu_{1i} = \beta_i \theta_{1i} + \gamma_3, \quad i \in \{1,2\} \label{eq:foc_z1} \\
\frac{\partial \mathcal{L}}{\partial q_i} &= -\beta_i \theta_{0i} + \gamma_i z_{1i} \frac{\partial \mu_{1i}}{\partial q_i} \quad  \nonumber \\
&\quad \text{Note: Treating } \mu_{1i} \text{ as state constraint } C q_i^{\alpha_1} z_0^{\alpha_2} - \mu_{1i} = 0 \text{ yields equivalent dual relations.} \nonumber \\
&\quad \text{Standard form derivation yields: } \beta_i \theta_{0i} = \gamma_i z_{1i}\alpha_1 \frac{\mu_{1i} }{q_i} = \gamma_i \alpha_1 \frac{\mu_{1i} z_{1i}}{q_i}, \label{eq:foc_q} \\
\frac{\partial \mathcal{L}}{\partial z_0} &= -\gamma_3 + \sum_{i=1}^2 \gamma_i z_{1i} \frac{\partial \mu_{1i}}{\partial z_0}= -\gamma_3 +  \sum_{i=1}^2 \gamma_i \alpha_2 \frac{\mu_{1i} z_{1i}}{z_0} = 0 \implies \gamma_3 = \alpha_2 \sum_{i=1}^2 \gamma_i \frac{\mu_{1i} z_{1i}}{z_0}. \label{eq:foc_z0}
\end{align}

\subsubsection{Derivation of the Optimality Conditions}
We now manipulate the FOCs to derive the key matching index (KMI) conditions stated in Theorem \ref{thm:general_optimality}.

\textbf{Deriving the Marginal Value Conservation (Proposition \ref{prop:first_level})}

We begin by substituting the expression for $\beta_i$ derived in \eqref{eq:foc_q} into \eqref{eq:foc_z1}.
Recalling from \eqref{eq:foc_q} that $\beta_i = \frac{\gamma_i}{\theta_{0i}} \alpha_1 \frac{\mu_{1i} z_{1i}}{q_i}$ (assuming unit consistency in differentials), the substitution into \eqref{eq:foc_z1} yields the relation $\gamma_i\mu_{1i}=  \left( \gamma_i \alpha_1 \frac{\mu_{1i}z_{1i}}{q_i\theta_{0i}} \right)\theta_{1i} + \gamma_3$.
Rearranging this equation to isolate $\gamma_3$ reveals that $\gamma_3  = \gamma_i\mu_{1i} \left( 1 - \alpha_1 \frac{z_{1i} \theta_{1i}}{q_i \theta_{0i}} \right)$.
By defining the elasticity term $s_i \triangleq 1 - \alpha_1 \frac{z_{1i} \theta_{1i}}{q_i \theta_{0i}}$, we simplify the relationship to
\begin{equation}
    \gamma_3 = \gamma_i \mu_{1i} s_i. \label{eq:gamma3_relation}
\end{equation}
Substituting \eqref{eq:gamma3_relation} into the driver conservation condition \eqref{eq:foc_z0} leads to the expression $\gamma_3 = \alpha_2 \left( \frac{\gamma_3}{\mu_{11} s_1} \frac{\mu_{11} z_{11}}{z_0} + \frac{\gamma_3}{\mu_{12} s_2} \frac{\mu_{12} z_{12}}{z_0} \right)$.
Since the shadow price of driver conservation $\gamma_3$ is non-zero, dividing through by $\gamma_3$ yields
\begin{equation}
    1 = \frac{\alpha_2 z_{11}/z_0}{s_1} + \frac{\alpha_2 z_{12}/z_0}{s_2},
\end{equation}
which recovers Eq. \eqref{eq:first_level}, representing the conservation of the marginal value of idle drivers.

\textbf{Deriving the KMI Unity Conditions}

We subsequently distinguish between the two objective cases to establish the KMI unity conditions.
For the case of Societal Welfare where $w_1=w_2=1$, the condition derived from \eqref{eq:foc_z2} implies $\gamma_3 = 1 - (\beta_i + \gamma_i)\mu_2$.
Substituting $\beta_i$ from \eqref{eq:foc_q} into this expression establishes a direct relationship between $\gamma_1$ and $\gamma_2$.
When combined with \eqref{eq:gamma3_relation}, this yields the condition for $\zeta_i^W = 1$ as defined in Definition 2.
The cross-term in $\zeta_i^W$ arises directly from the term $\gamma_j$ appearing in the summation of \eqref{eq:foc_z0}, representing the opportunity cost of allocating an idle driver to type $j$ instead of $i$.

Alternatively, for the case of Platform Revenue where $w_i = p_i$, the condition from \eqref{eq:foc_z2} becomes $\gamma_3 = p_i - (\beta_i + \gamma_i)\mu_2$.
The presence of $p_i$ introduces asymmetry into the system.
Letting $A_i = 1 + \alpha_1 \frac{\mu_{1i} z_{1i}}{q_i \theta_{0i}}$, the multiplier relationship transforms into
\begin{equation}
    \gamma_i = \frac{p_i}{\mu_{1i} s_i + \mu_2 A_i}.
\end{equation}
Substituting this result into the conservation law \eqref{eq:foc_z0} and rearranging terms yields the Revenue KMI expression $\zeta_i^R$.
Here, the ratio $p_j/p_i$ and the terms $A_i, A_j$ in $\zeta_i^R$ explicitly capture the substitution of value (revenue) rather than volume (matches).

\subsubsection{Remark on Sufficiency and Global Optimality}
\label{remark:sufficiency}
The conditions derived in Theorem \ref{thm:general_optimality} characterize a KKT stationary point $(\boldsymbol{\mu}^*, \mathbf{x}^*)$.
In general nonlinear optimization, stationarity is a necessary but not sufficient condition for global optimality, particularly given the non-convex nature of the spatial matching function $\mu_i(z_0, q_i) = C q_i^{\alpha_1} z_0^{\alpha_2}$.

While the KMI condition $\boldsymbol{\zeta}(\boldsymbol{\mu}) = \mathbf{1}$ formally characterizes a stationary point, the structural properties of the system suggest this solution corresponds to the global optimum in the relevant operational regime.
Central to this argument is the uniqueness of the fluid equilibrium $x^*$ for any fixed set of thresholds $\boldsymbol{\mu}$, as established in Theorem~\ref{thm:EC4}, which effectively precludes the emergence of multiple steady states for a single policy.
This theoretical foundation is reinforced by extensive numerical evidence, visualized in Figures~\ref{fig:3d_objectives} and~\ref{fig:3d_kmi}, indicating that the objective function $J(\boldsymbol{\mu})$ behaves unimodally, specifically displaying quasi-concavity, with respect to $\boldsymbol{\mu}$ in the region where $z_0 > 0$.
Furthermore, the Jacobian $\mathbf{J}_\zeta$ maintains full rank and diagonal dominance, a characteristic that ensures the gradient-based control law follows a unique ascent path.
Consequently, despite the formal limitation to stationarity, the convergence of theoretical uniqueness and empirical unimodality strongly supports the conclusion that the solution is the global maximizer for the platform's objective.

\subsection{Structural Divergence of Welfare and Revenue Indices}
\label{sec:ec_kmi_difference}

This section shows why the Revenue KMI $\zeta_i^R$ does not algebraically reduce to the Welfare KMI $\zeta_i^W$ even in the symmetric price case ($p_1=p_2=1$).
This divergence is not an artifact of derivation but a reflection of structural economic differences in how opportunity costs are valued.

The distinction arises primarily from the contrast between physical and economic opportunity costs.
While $\zeta_i^W$ optimizes the aggregate rate of service ($z_{21} + z_{22}$), the Lagrange multipliers associated with this objective represent the shadow price of system constraints in units of throughput per hour.
Consequently, the cross-market term in $\zeta_i^W$ penalizes the reduction in matching volume for class $j$ caused by serving class $i$.
In contrast, $\zeta_i^R$ optimizes accumulated value ($p_1 z_{21} + p_2 z_{22}$), effectively treating each match as a unit of revenue generation rather than a raw throughput event.
Even when prices are symmetric, the multipliers $\gamma_i$ in the revenue case incorporate the price $p_i$ directly into the numerator, such that $\gamma_i \propto \frac{p_i}{\mu_{1i} s_i + \mu_2 A_i}$, thereby embedding the marginal revenue substitution rate directly into the KMI.

Mathematically, this structural difference is evident in the multiplier relationships derived in the proof.
For the welfare objective, the relation between $\gamma_1$ and $\gamma_2$ depends purely on the ratio of service completion rates and queue elasticities.
Conversely, for the revenue objective, the relation is governed by $\frac{\gamma_1}{\gamma_2} = \frac{p_1 (\mu_{12} s_2 + \mu_2 A_2)}{p_2 (\mu_{11} s_1 + \mu_2 A_1)}$.
The terms $A_i = 1 + \alpha_1 \frac{\mu_{1i} z_{1i}}{q_i \theta_{0i}}$ appear specifically in the revenue formulation because the cost of a driver is weighed against the revenue $p_i$, which is coupled with the trip duration and turnover rate.
This creates a structurally distinct denominator in $\zeta_i^R$ that persists regardless of whether $p_i = p_j$.
Therefore, the two objectives imply distinct operational frontiers; optimizing for revenue, even with equal prices, requires a different balancing of queue accumulation compared to maximizing raw throughput, as the platform implicitly values the reliability of revenue generation differently from the volume of passenger flow.

\section{Proof of Stochastic Model Converges to Fluid Model} 
\label{sec:ec_stochastic_to_fluid}

% Define custom command for positive part of ceiling
\newcommand{\ceilpos}[1]{\left\lceil #1 \right\rceil_+}
\newcommand{\plus}{_+}

We establish the convergence of stochastic model to fluid model 
when $\alpha_1 = \alpha_2 = \alpha$. Recall that we apply a threshold rule 
to control the matching process, i.e., a matching happens if and only 
if $\mu_{11}(t) \geq \mu_{11}$ or $\mu_{12}(t) \geq \mu_{12}$ for some thresholds value $\mu_{11}, \mu_{12}$. 
By Assumption~\ref{eq:spatial_model_rate} , when $\alpha_1 = \alpha_2$, a matching happens if and only if 
$Z_0(t)Q_i(t) \ge \bar{\mu}_{1i}$, where $\bar{\mu}_{1i} =\sqrt[\alpha]{\frac{\mu_{1i}}{C}}$ for some $\mu_{1i}, i = 1,2$.
When $Z_0(t)Q_i(t) \le \bar{\mu}_{1i}$, the system dynamics of the stochastic model are described by equations (4)--(7), and the matching equations (9)--(10) can be rewritten in the following equivalent form:
\begin{equation}
\begin{aligned}
M(t) = M_1(t) + M_2(t) &= \int_0^t \mathbf{1}_{\{(Q_1(s-)+1)Z_0(s-) \ge \bar{\mu}_{11}\}}d A_1(s) 
+ \int_0^t \mathbf{1}_{\{(Q_2(s-)+1)Z_0(s-) \ge \bar{\mu}_{12}\}} d A_2(s)\\
&+ \int_0^t  \mathbf{1}_{\{Q_1(s-)(Z_0(s-)+1) \ge \bar{\mu}_{11} \ \text{or} \ Q_2(s-)(Z_0(s-)+1)  \ge \bar{\mu}_{12}\}} d( \sum_{i=1}^2 (R_{1i}(s) + D_{2i}(s))), 
\end{aligned}
\label{eq:EC33}
\end{equation}
where the first two terms correspond to a matching triggered by a passenger arrival and the third term corresponds to a matching triggered by a driver becoming idle.

Before presenting the detailed proof, we provide a roadmap. 
As a first step, we need to characterize the state evolution dynamics of the stochastic 
model using the system input processes and parameters. 
The main difficulty lies in the nontrivial dependence of the matching process $M(t)$ 
on the system inputs, which is tackled by Lemma~\ref{lem:EC1}. Having established the dynamic equations, 
we observe that the stochastic and fluid model both satisfy differential equations of form 
\eqref{eq:EC39}-\eqref{eq:EC42}, with difference in residual terms $B_i(t), i=1,2,3,4$. 
For given $B_i(t), i = 1,2,3,4$, we first show that \eqref{eq:EC39}-\eqref{eq:EC42} can 
be regarded as a contraction mapping (Proposition~\ref{prop:EC1}), thus validating the existence and 
uniqueness of the stochastic model. Next, we establish using the Law of Large Numbers that 
the difference in residual terms $B_i(t), i = 1,2,3,4$, is asymptotically negligible (Lemma~\ref{lem:EC2}). 
Thus, combined with the fact that \eqref{eq:EC39}-\eqref{eq:EC42} is Lipschitz continuous, 
we complete the proof of convergence (Theorem~\ref{thm:EC1}).

We begin with Lemma~\ref{lem:EC1}, which characterizes the matching process under the threshold control policy; the proof proceeds by induction.
Here $\lceil x \rceil$ denotes the smallest integer greater than or equal to $x$.

\begin{lemma} \label{lem:EC1}
Let $X_i(t) = Q_i(0) + A_i(t) - R_{0i}(t)$, $X(t)=X_1(t) + X_2(t)$ and $Y(t) = Z_0(0) + R_{11}(t) + D_{21}(t)+ R_{12}(t) + D_{22}(t)$, then
\begin{equation}
M(t) = M_1(t)+M_2(t)
= \sup_{0 \le s \le t} \ceilpos{\frac{X(s) + Y(s) - \sqrt{(X(s) - Y(s))^2 + 4(\bar{\mu}_{11}+\bar{\mu}_{12})}}{2}} . \label{eq:EC34}
\end{equation}
\end{lemma}

\begin{proof}{Proof of Lemma~\ref{lem:EC1}}
The proof uses sample-path arguments.
First we introduce the functions of calculating the amount of the matched two different class passengers.
Let $x_1$ and $x_2$ be the number of matched passengers of class 1 and class 2, respectively, 
and $x = x_1 + x_2$, the total number of matched passengers.
We can write the matching process as follows:
\begin{equation*}
\begin{split}
    (X_1(t) - x_1)(Y(t) - x_1-x_2) = \bar{\mu}_{11},\\
    (X_2(t) - x_2)(Y(t) - x_1-x_2) = \bar{\mu}_{12}.\\
\end{split}
\end{equation*}
Adding these two equations yields
\begin{equation*}
\begin{split}
    (X_1(t) - x_1)(Y(t) - x_1-x_2) + (X_2(t) - x_2)(Y(t) - x_1-x_2)\\
    = (X(t) - x)(Y(t) - x) = \bar{\mu}_{11} + \bar{\mu}_{12}.
\end{split}
\end{equation*}

Observe that $\frac{X(s) + Y(s) - \sqrt{(X(s) - Y(s))^2 + 4(\bar{\mu}_{11}+\bar{\mu}_{12})}}{2}$ is the root of the equation $(X(t) - x)(Y(t) - x) = \bar{\mu}_{11} + \bar{\mu}_{12}$.
Assume that the $k$-th matching happens at 
time $t_k$ with $0 < t_1 < t_2 < \dots < t_k < \dots$ . For any $t \in [0, t_1)$, 
we have $X_1(t) = Q_1(t),X_2(t) = Q_2(t), Y(t) = Z_0(t),$ and $X_i(t)Y(t) \le \bar{\mu}_{1i}, i=1,2,$ which implies 
that the $X(t)Y(t) \le \bar{\mu}_{11}+\bar{\mu}_{12},$ so the right-hand side of \eqref{eq:EC34} is 0. When $t = t_1$, since the first 
matching happened, we have $X(t_1) = Q_1(t_1)+Q_2(t_1) + 1, Y(t_1) = Z_0(t_1) + 1, X(t_1)Y(t_1) > \bar{\mu}_{11}+\bar{\mu}_{12},$ 
and $(X(t_1) - 1)(Y(t_1) - 1) \le \bar{\mu}_{11}+\bar{\mu}_{12}.$
Since $\sup_{0 \le s \le t} \frac{X(s)+Y(s) - \sqrt{(X(s)-Y(s))^2+4(\bar{\mu}_{11}+\bar{\mu}_{12})}}{2}$ 
is the root to equation $(X(t) - x)(Y(t) - x) = \bar{\mu}_{11}+\bar{\mu}_{12},$ the right-hand side 
of \eqref{eq:EC34} is 1 when $t=t_1$. By induction, one can verify that 
when $t \in [t_k, t_{k+1})$, $X(t) - k = Q_1(t)+Q_2(t), Y(t) - k = Z_0(t),$ $(X(t) - k+1)(Y(t) - k+1) > \bar{\mu}_{11}+\bar{\mu}_{12},$ 
and $(X(t)-k)(Y(t)-k) \le \bar{\mu}_{11}+\bar{\mu}_{12}$. These imply that the right-hand side of 
\eqref{eq:EC34} is $k$ for $t \in [t_k, t_{k+1})$ and prove the result. 
\Halmos
\end{proof}

To carry out our asymptotic analysis, we index the performance of the sequence of queues by a superscript $n$, denoted by $Q_i^n, Z^n_0, Z^n_{1i}, Z^n_{2i}, M_i^n, R^n_{0i}, R^n_{1i}, D^n_{1i}, D^n_{2i}, X_i^n, i=1,2,$ and $Y^n$, where the scaling is defined in the beginning of Section~\ref{sec:Fluid_Model}.
According to \eqref{eq:EC34}, for the $n$-th system, the system dynamics can be characterized as follows,

\begin{align}
i = 1,2: \quad
Q_i^n(t) &= Q_i^n(0) + A_i^n(t) - R_{0i}^n(t) - M_i^n(t), \label{eq:EC35} \\
Z^n_0(t) &= Z^n_0(0) + \sum_{i=1}^{2} \{R^n_{1i}(t) + D^n_{2i}(t) - M_i^n(t)\}, \label{eq:EC36} \\
Z^n_{1i}(t) &= Z^n_{1i}(0) + M_i^n(t) - D^n_{1i}(t) - R^n_{1i}(t), \label{eq:EC37} \\
Z^n_{2i}(t) &= Z^n_{2i}(0) + D^n_{1i}(t) - D^n_{2i}(t). \label{eq:EC38}
\end{align}
We add a bar to denote quantities of fluid-scaled processes, e.g., $\bar{Q}_1^n(t) = Q_1^n(t)/n$.
By using centering operation, the fluid-scaled dynamics can be represented as follows
\begin{align}
    i = 1,2: \quad
\bar{Q}_i^n(t) &= (\bar{X}^n_i(t)-\tilde{X}_i^n(t)) - (\bar{M}_i^n(t) - \tilde{M}_i^n(t)) + \tilde{X}_i^n(t) - \tilde{M}_i^n(t), \label{eq:EC39b} \\
\bar{Z}^n_0(t) &= (\bar{Y}^n(t) - \tilde{Y}^n(t)) - \sum_{i=1}^{2}(\bar{M}_i^n(t) - \tilde{M}_i^n(t)) + \tilde{Y}^n(t) - \sum_{i=1}^{2}\tilde{M}_i^n(t), \label{eq:EC40b} \\
\bar{Z}^n_{1i}(t) &= \bar{Z}^n_{1i}(0) - (\bar{R}^n_{1i}(t) - \theta_{1i} \int_0^t \bar{Z}^n_{1i}(s)ds) - (\bar{D}^n_{1i}(t) - \mu_{1i} \int_0^t \bar{Z}^n_{1i}(s)ds) \nonumber \\
& \quad + (\bar{M}^n_{i}(t) - \tilde{M}^n_{i}(t)) + \tilde{M}^n_{i}(t) - \theta_{1i} \int_0^t \bar{Z}^n_{1i}(s)ds - \mu_{1i} \int_0^t \bar{Z}^n_{1i}(s)ds, \label{eq:EC41b} \\
\bar{Z}^n_{2i}(t) &= \bar{Z}^n_{2i}(0) + (\bar{D}^n_{1i}(t) - \mu_{1i} \int_0^t \bar{Z}^n_{1i}(s)ds) - (\bar{D}^n_{2i}(t) - \mu_2 \int_0^t \bar{Z}^n_{2i}(s)ds) \nonumber \\
& \quad + \mu_{1i} \int_0^t \bar{Z}^n_{1i}(s)ds - \mu_2 \int_0^t \bar{Z}^n_{2i}(s)ds, \label{eq:EC42b}
\end{align}
where $\tilde{M}^n(t)= \tilde{M}_1^n(t)+\tilde{M}_2^n(t)= \sup_{0 \le s \le t} \ceilpos{\frac{\tilde{X}^n(s) + \tilde{Y}^n(s) - \sqrt{(\tilde{X}^n(s) - \tilde{Y}^n(s))^2 + 4(\bar{\mu}_{11}+\bar{\mu}_{12})}}{2}},$
$\tilde{X}^n_i(t) = \bar{Q}_i^n(0) + \lambda_i t - \theta_{0i} \int_0^t \bar{Q}_i^n(s)ds,i =1,2$, and $\tilde{X}^n(t)=\tilde{X}^n_1(t)+\tilde{X}^n_2(t)$, $\tilde{Y}^n(t) = \bar{Z}^n_0(t) + \sum_{i=1}^{2}(\theta_{1i} \int_0^t \bar{Z}^n_{1i}(s)ds + \mu_2 \int_0^t \bar{Z}^n_{2i}(s)ds).$
Note that $\tilde{M}^n(t)$ is the centering process of the matching process $\bar{M}^n(t)$.
We have the following result.

\begin{proposition} \label{prop:EC1}
For any given $(B_{1i}, B_2, B_{3i}, B_{4i}, i =1,2) \in D([0, T], \mathbb{R}^7)$, there exists a unique 
$(Q_i, Z_0, Z_{1i}, Z_{2i}, i=1,2) \in D([0, T], \mathbb{R}^7)$ satisfying
\begin{align}
    i = 1,2: \quad
Q_i(t) &= \tilde{X}_i(t) + B_{1i}(t) - \tilde{M}_{i}(t), \label{eq:EC39} \\
Z_0(t) &= \tilde{Y}(t) + B_2(t) - \tilde{M}_{1}(t)-\tilde{M}_{2}(t), \label{eq:EC40} \\
Z_{1i}(t) &= Z_{1i}(0) + B_{3i}(t) - \theta_{1i} \int_0^t Z_{1i}(s)ds - \mu_{1i} \int_0^t Z_{1i}(s)ds + \tilde{M}_{i}(t), \label{eq:EC41} \\
Z_{2i}(t) &= Z_{2i}(0) + B_{4i}(t) + \mu_{1i} \int_0^t Z_{1i}(s)ds - \mu_2 \int_0^t Z_{2i}(s)ds, \label{eq:EC42}
\end{align}
where $\tilde{M}(t)= \tilde{M}_1(t)+\tilde{M}_2(t) = \sup_{0 \le s \le t} \lceil{\frac{\tilde{X}(s) + \tilde{Y}(s) - \sqrt{(\tilde{X}(s) - \tilde{Y}(s))^2 + 4(\bar{\mu}_{11}+\bar{\mu}_{12})}}{2}}\rceil^+,$
$\tilde{X}_i(t) = Q_i(0) + \lambda_i t - \theta_{0i} \int_0^t Q_i(s)ds, i=1,2,$ $\tilde{X}(s) = \tilde{X}_1(s)+\tilde{X}_2(s)$,
and $\tilde{Y}(t) = Z_0(t) + \sum_{i=1}^{2}(\theta_{1i} \int_0^t Z_{1i}(s)ds + \mu_2 \int_0^t Z_{2i}(s)ds).$ 
Thus, these representations constitute a mapping $f$ from $(B_{1i}, B_2, B_{3i}, B_{4i},i=1,2) \in D([0, T], \mathbb{R}^7)$ 
to $(Q_i, Z_0, Z_{1i}, Z_{2i},i=1,2) \in D([0, T], \mathbb{R}^7)$. In addition, the mapping $f$ is Lipschitz continuous under the uniform norm.
\end{proposition}

\begin{proof}{Proof of Proposition~\ref{prop:EC1}}
The proof is based on the contraction mapping principle.
For any $(Q_i, Z_0, Z_{1i}, Z_{2i},i=1,2) \in D([0, T], \mathbb{R}^7)$ and $b < \frac{1}{\sum_{i=1}^{2}(4\theta_{0i} + 5\theta_{1i} + 2\mu_{1i} + 5\mu_2)}$, define a mapping $\Phi$ from $D([0, b], \mathbb{R}^4)$ to $D([0, b], \mathbb{R}^7)$ given by \eqref{eq:EC39}-\eqref{eq:EC42}.
We prove that $\Phi$ is a contraction mapping.
For any $(Q_i, Z_0, Z_{1i}, Z_{2i},i=1,2) \in D([0, T], \mathbb{R}^7)$ and $(Q_i^*, Z^*_0, Z^*_{1i}, Z^*_{2i},i=1,2) \in D([0, b], \mathbb{R}^7)$, we have

\begin{align}
& \| \Phi(Q_i, Z_0, Z_{1i}, Z_{2i},i=1,2) - \Phi(Q_i^*, Z^*_0, Z^*_{1i}, Z^*_{2i},i=1,2) \|_b \nonumber\\
& \le \sum_{i=1}^{2}\{\theta_{0i} \int_0^b |Q_i(s) - Q_i^*(s)|ds + 2(\theta_{1i} + \mu_{1i}) \int_0^b |Z_{1i}(s) - Z^*_{1i}(s)|ds + 2\mu_2 \int_0^b |Z_{2i}(s) - Z^*_{2i}(s)|ds\} \nonumber\\
& + 3 \sup_{0 \le s \le b} |\tilde{M}(s) - \tilde{M}^*(s)|. \label{eq:EC16}
\end{align}

For any $s \in [0, b]$, we have
\begin{equation*}
\begin{aligned}
& |\tilde{M}(s) - \tilde{M}^*(s)| \\
& \leq \sup_{0 \leq r \leq s} \left| \left\lceil \frac{\tilde{X}(r) + \tilde{Y}(r) - \sqrt{(\tilde{X}(r) - \tilde{Y}(r))^2 + 4(\bar{\mu}_{11}+\bar{\mu}_{12})}}{2} \right\rceil^+ \right. \\
& \quad\quad\quad\quad\quad\quad - \left. \left\lceil \frac{\tilde{X}^*(r) + \tilde{Y}^*(r) - \sqrt{(\tilde{X}^*(r) - \tilde{Y}^*(r))^2 + 4(\bar{\mu}_{11}+\bar{\mu}_{12})}}{2} \right\rceil^+ \right| \\
& \leq \sup_{0 \leq r \leq s} \left| \frac{\tilde{X}(r) + \tilde{Y}(r) - \sqrt{(\tilde{X}(r) - \tilde{Y}(r))^2 + 4(\bar{\mu}_{11}+\bar{\mu}_{12})}}{2} \right. \\
& \quad\quad\quad\quad\quad\quad - \left. \frac{\tilde{X}^*(r) + \tilde{Y}^*(r) - \sqrt{(\tilde{X}^*(r) - \tilde{Y}^*(r))^2 + 4(\bar{\mu}_{11}+\bar{\mu}_{12})}}{2} \right| \\
& \leq \sup_{0 \leq r \leq s} \left| \frac{\tilde{X}(r) - \tilde{X}^*(r) + \tilde{Y}(r) - \tilde{Y}^*(r)}{2} \right. \\
& \quad\quad\quad\quad\quad\quad + \left. \frac{[\tilde{X}(r) - \tilde{Y}(r) - (\tilde{X}^*(r) - \tilde{Y}^*(r))][\tilde{X}(r) - \tilde{Y}(r) + \tilde{X}^*(r) - \tilde{Y}^*(r)]}{2[\sqrt{(\tilde{X}(r) - \tilde{Y}(r))^2 + 4(\bar{\mu}_{11}+\bar{\mu}_{12})} + \sqrt{(\tilde{X}^*(r) - \tilde{Y}^*(r))^2 + 4(\bar{\mu}_{11}+\bar{\mu}_{12})}]} \right| \\
& \leq \sup_{0 \leq r \leq s} \left( |\tilde{X}(r) - \tilde{X}^*(r)| + |\tilde{Y}(r) - \tilde{Y}^*(r)| \right) \\
& \leq \sup_{0 \leq r \leq s} \left( |\tilde{X}_1(r) - \tilde{X}_1^*(r)| + |\tilde{X}_2(r) - \tilde{X}_2^*(r)| + |\tilde{Y}(r) - \tilde{Y}^*(r)| \right) \\
& \leq \sum_{i=1}^{2}(\theta_{0i} \int_0^s |Q_i(r) - Q_i^*(r)| \, dr + \theta_{1i} \int_0^s |Z_{1i}(r) - Z_{1i}^*(r)| \, dr + \mu_2 \int_0^s |Z_{2i}(r) - Z_{2i}^*(r)| \, dr )\\
& \leq \sum_{i=1}^{2}(\theta_{0i} \int_0^b |Q_i(r) - Q_i^*(r)| \, dr + \theta_{1i} \int_0^b |Z_{1i}(r) - Z_{1i}^*(r)| \, dr + \mu_2 \int_0^b |Z_{2i}(r) - Z_{2i}^*(r)| \, dr).
\end{aligned}
\end{equation*}

Thus we have
\begin{align}
\sup_{0 \le s \le b} |\tilde{M}(s) - \tilde{M}^*(s)| &\le \sum_{i=1}^{2}\left(\theta_{0i} \int_0^b |Q_i(r) - Q_i^*(r)| \, dr \right. \nonumber \\
&\quad + \theta_{1i} \int_0^b |Z_{1i}(r) - Z_{1i}^*(r)| \, dr \nonumber \\
&\quad \left. + \mu_2 \int_0^b |Z_{2i}(r) - Z_{2i}^*(r)| \, dr\right). \label{eq:EC43}
\end{align}

Substituting \eqref{eq:EC43} into \eqref{eq:EC16} yields
\begin{align}
&\| \Phi(Q_i, Z_0, Z_{1i}, Z_{2i}, i=1,2) - \Phi(Q_i^*, Z^*_0, Z^*_{1i}, Z^*_{2i}, i=1,2) \|_b \nonumber \\
&\le \sum_{i=1}^{2}\left\{\theta_{0i} \int_0^b |Q_i(s) - Q_i^*(s)|ds + 2(\theta_{1i} + \mu_{1i}) \int_0^b |Z_{1i}(s) - Z^*_{1i}(s)|ds \right. \nonumber \\
&\quad\quad\quad\quad\quad + \left. 2\mu_2 \int_0^b |Z_{2i}(s) - Z^*_{2i}(s)|ds\right\} \nonumber \\
&\quad + 3 \sum_{i=1}^{2}\left(\theta_{0i} \int_0^b |Q_i(r) - Q_i^*(r)| \, dr + \theta_{1i} \int_0^b |Z_{1i}(r) - Z_{1i}^*(r)| \, dr \right. \nonumber \\
&\quad\quad\quad\quad\quad + \left. \mu_2 \int_0^b |Z_{2i}(r) - Z_{2i}^*(r)| \, dr\right) \nonumber \\
&= \sum_{i=1}^{2}\left\{(4\theta_{0i}) \int_0^b |Q_i(s) - Q_i^*(s)|ds + (5\theta_{1i} + 2\mu_{1i}) \int_0^b |Z_{1i}(s) - Z^*_{1i}(s)|ds \right. \nonumber \\
&\quad\quad\quad\quad\quad + \left. (5\mu_2) \int_0^b |Z_{2i}(s) - Z^*_{2i}(s)|ds\right\} \nonumber \\
&\le \left(\sum_{i=1}^{2}(4\theta_{0i} + 5\theta_{1i} + 2\mu_{1i} + 5\mu_2)\right)b \|(Q_i, Z_0, Z_{1i}, Z_{2i}, i=1,2) - (Q_i^*, Z^*_0, Z^*_{1i}, Z^*_{2i}, i=1,2)\|_b \nonumber \\
&< \|(Q_i, Z_0, Z_{1i}, Z_{2i}, i=1,2) - (Q_i^*, Z^*_0, Z^*_{1i}, Z^*_{2i}, i=1,2)\|_b.
\end{align}

Thus, the mapping $\Phi$ is a contraction. By the Contraction Mapping Theorem, there exists a unique solution to \eqref{eq:EC39}--\eqref{eq:EC42} on $[0, b]$. It now remains to extend the existence and uniqueness result from $[0, b]$ to $[0, T]$. 

To extend to the full interval $[0, T]$, we proceed iteratively. Replacing $(Q_i(t), Z_0(t), Z_{1i}(t), Z_{2i}(t))$ by $(Q_i(b+t), Z_0(b+t), Z_{1i}(b+t), Z_{2i}(b+t))$ for $i = 1,2$, and following the previous argument, we obtain a unique solution for \eqref{eq:EC39}--\eqref{eq:EC42} on $[0, 2b]$. Repeating this approach for $\lceil T/b \rceil$ times gives a unique solution on the interval $[0, T]$.

Next we prove that the mapping $f$ is Lipschitz continuous under the uniform norm.
For any
\begin{equation*}
  \|(B_{1i}, B_2, B_{3i}, B_{4i}, i=1,2) - (B^*_{1i}, B^*_2, B^*_{3i}, B^*_{4i}, i=1,2)\|_T < \epsilon,
\end{equation*}
following the same steps as in the proof of Proposition~\ref{prop:EC1}, we obtain
\begin{align}
&\|f(B_{1i}, B_2, B_{3i}, B_{4i}, i=1,2) - f(B^*_{1i}, B^*_2, B^*_{3i}, B^*_{4i}, i=1,2)\|_T \nonumber \\
&\leq 4\|(B_{1i}, B_2, B_{3i}, B_{4i}, i=1,2) - (B^*_{1i}, B^*_2, B^*_{3i}, B^*_{4i}, i=1,2)\|_T \nonumber \\
&\quad + \sum_{i=1}^{2}\left\{4\theta_{0i} \int_0^T |Q_i(s) - Q_i^*(s)|ds + (5\theta_{1i} + 2\mu_{1i}) \int_0^T |Z_{1i}(s) - Z^*_{1i}(s)|ds \right. \nonumber \\
&\quad\quad\quad\quad\quad + \left. 5\mu_2 \int_0^T |Z_{2i}(s) - Z^*_{2i}(s)|ds\right\} \nonumber \\
&\leq 4\epsilon + \left(\sum_{i=1}^{2}(4\theta_{0i} + 5\theta_{1i} + 2\mu_{1i} + 5\mu_2)\right) \nonumber \\
&\quad \times \int_0^T \|f(B_{1i}, B_2, B_{3i}, B_{4i}, i=1,2) - f(B^*_{1i}, B^*_2, B^*_{3i}, B^*_{4i}, i=1,2)\|_s \, ds.
\end{align}

Thus by Grönwall's inequality, we have
\begin{align}
&\|f(B_{1i}, B_2, B_{3i}, B_{4i}, i=1,2) - f(B^*_{1i}, B^*_2, B^*_{3i}, B^*_{4i}, i=1,2)\|_T \nonumber \\
&\leq 4\epsilon \exp\left\{\left(\sum_{i=1}^{2}(4\theta_{0i} + 5\theta_{1i} + 2\mu_{1i} + 5\mu_2)\right)T\right\},
\end{align}
which implies that the mapping $f$ is Lipschitz continuous under the uniform norm. 
\Halmos
\end{proof}

When $n \to \infty$, the weak convergence of the centered processes
\(
\bar{A}_i^n(t) - \lambda_i t, \quad \bar{R}^n_{0i}(t) - \theta_{0i} \int_0^t \bar{Q}_i^n(s)ds, \quad \bar{R}^n_{1i}(t) - \theta_{1i} \int_0^t \bar{Z}^n_{1i}(s)ds,
\)
and
\(
\bar{D}^n_{2i}(t) - \mu_2 \int_0^t \bar{Z}^n_{2i}(s)ds \quad \text{for } i = 1,2
\)
follows from the functional strong law of large numbers. We further prove that $\bar{M}^n(t) - \tilde{M}^n(t)$ and $\bar{D}^n_{1i}(t) - \mu_{1i} \int_0^t \bar{Z}^n_{1i}(s)ds$ converge to 0 when $n \to \infty$ for $i = 1,2$.

\begin{lemma} \label{lem:EC2}
For any $\epsilon > 0$ and any given $T$, we have
\begin{equation}
\lim_{n \to \infty} \mathbb{P} \left( \sup_{t \le T} |\bar{M}^n(t) - \tilde{M}^n(t)| > \epsilon \right) = 0, \label{eq:EC44}
\end{equation}
and
\begin{equation}
\lim_{n \to \infty} \mathbb{P} \left( \sup_{t \le T} \left|\bar{D}^n_{1i}(t) - \mu_{1i} \int_0^t \bar{Z}^n_{1i}(s)ds\right| > \epsilon \right) = 0, \quad i = 1,2. \label{eq:EC45}
\end{equation}
\end{lemma}
\begin{proof}{Proof of Lemma~\ref{lem:EC2}}
We prove each result separately.

\textbf{Part 1: Convergence of matching process}

For any $t \le T$, by Lemma~\ref{lem:EC1} and the definition of $\tilde{M}^n(t)$, we have
\begin{align}
|\bar{M}^n(t) - \tilde{M}^n(t)| &\le \frac{1}{n} + \left| \sup_{0 \le s \le t} \left\lceil \frac{\bar{X}^n(s) + \bar{Y}^n(s) - \sqrt{(\bar{X}^n(s) - \bar{Y}^n(s))^2 + 4(\bar{\mu}_{11}+\bar{\mu}_{12})}}{2} \right\rceil^+ \right. \nonumber \\
&\quad - \left. \sup_{0 \le s \le t} \left\lceil \frac{\tilde{X}^n(s) + \tilde{Y}^n(s) - \sqrt{(\tilde{X}^n(s) - \tilde{Y}^n(s))^2 + 4(\bar{\mu}_{11}+\bar{\mu}_{12})}}{2} \right\rceil^+ \right| \nonumber \\
&\le \frac{1}{n} + \sup_{0 \le s \le t} |\bar{X}^n(s) - \tilde{X}^n(s)| + \sup_{0 \le s \le t} |\bar{Y}^n(s) - \tilde{Y}^n(s)|.
\end{align}

Now we bound each term:
\begin{align}
\sup_{0 \le s \le t} |\bar{X}^n(s) - \tilde{X}^n(s)| &= \sup_{0 \le s \le t} |\bar{X}^n_1(s) + \bar{X}^n_2(s) - \tilde{X}^n_1(s) - \tilde{X}^n_2(s)| \nonumber \\
&\le \sup_{0 \le s \le t} \sum_{i=1}^{2} \left| (\bar{A}^n_i(s) - \lambda_i s) - (\bar{R}^n_{0i}(s) - \theta_{0i} \int_0^s \bar{Q}^n_i(r)dr) \right|,
\end{align}
and
\begin{align}
\sup_{0 \le s \le t} |\bar{Y}^n(s) - \tilde{Y}^n(s)| &\le \sup_{0 \le s \le t} \sum_{i=1}^{2} \left| (\bar{R}^n_{1i}(s) - \theta_{1i} \int_0^s \bar{Z}^n_{1i}(r)dr) \right. \nonumber \\
&\quad + \left. (\bar{D}^n_{2i}(s) - \mu_2 \int_0^s \bar{Z}^n_{2i}(r)dr) \right|.
\end{align}

By the functional strong law of large numbers, each centered process converges to 0 in probability, which establishes \eqref{eq:EC44}.

\textbf{Part 2: Convergence of departure process}

For the convergence of $\bar{D}^n_{1i}(t) - \mu_{1i} \int_0^t \bar{Z}^n_{1i}(s)ds$, we first establish that the pickup rates remain bounded. Under our threshold control policy, a matching occurs when
\begin{equation}
\bar{Q}^n_i(t^-)\bar{Z}^n_0(t^-) \ge \bar{\mu}_{1i}, \quad i = 1,2.
\end{equation}

For any $\epsilon_1 > 0$, we show that
\begin{align}
&\lim_{n \to \infty} \mathbb{P} \left( \sup_{t \le T} \frac{(Q^n_i(t^-)+1)Z^n_0(t^-)}{n^2} \ge \bar{\mu}_{1i} + \epsilon_1 \right)+ \mathbb{P} \left( \sup_{t \le T} \frac{Q^n_i(t^-)(Z^n_0(t^-)+1)}{n^2} \ge \bar{\mu}_{1i} + \epsilon_1 \right) = 0. \label{eq:EC46}
\end{align}

The first term corresponds to a rider arrival triggering a match, while the second term corresponds to an empty car arrival triggering a match. 

Since $\frac{Q^n_i(t^-)}{n} \cdot \frac{Z^n_0(t^-)}{n} \le \bar{\mu}_{1i}$ by our control policy, the first case would require $Z^n_0(t^-) > \epsilon_1 n^2$. However, this contradicts the constraint $Z^n_0(t^-) \le n$ for sufficiently large $n$.

Similarly, the second case would require $Q^n_i(t^-) > \epsilon_1 n^2$. Since $Q^n_i(t^-) \le A^n_i(t^-) + Q^n_i(0)$, we would need $A^n_i(t^-) > \epsilon_1 n^2 - \bar{Q}_i(0)n$. By the law of large numbers:
\begin{equation}
\lim_{n \to \infty} \mathbb{P}\left(\sup_{t \le T} A^n_i(t^-) > \epsilon_1 n^2 - \bar{Q}_i(0)n\right) = 0.
\end{equation}

This establishes \eqref{eq:EC46}.

Given \eqref{eq:EC46}, we can apply a coupling argument: for large $n$, the departure process $D^n_{1i}(t)$ is sandwiched between two non-homogeneous Poisson processes with rates $\mu_{1i} Z^n_{1i}(t)$ and $(\mu_{1i} + \epsilon_1)Z^n_{1i}(t)$ respectively, with probability approaching 1. As $\epsilon_1 \to 0$, this yields the desired convergence in \eqref{eq:EC45}.
\Halmos
\end{proof}

\begin{theorem} \label{thm:EC1}
Suppose the initial condition holds. Then as $n \to \infty$, $(\bar{Q}^n_i, \bar{Z}^n_0, \bar{Z}^n_{1i}, \bar{Z}^n_{2i}, i=1,2)$ 
converges to $(\bar{Q}_i, \bar{Z}_0, \bar{Z}_{1i}, \bar{Z}_{2i}, i=1,2)$ in probability, where the limit satisfies 
the fluid model equations:
\begin{align*}
i = 1,2: \quad
\bar{Q}_i(t) &= \bar{Q}_i(0) + \lambda_i t - \theta_{0i} \int_0^t \bar{Q}_i(s)ds - \bar{M}_i(t), \\
\bar{Z}_0(t) &= \bar{Z}_0(0) + \sum_{i=1}^{2} \{\theta_{1i} \int_0^t \bar{Z}_{1i}(s)ds + \mu_2 \int_0^t \bar{Z}_{2i}(s)ds - \bar{M}_i(t)\}, \\
\bar{Z}_{1i}(t) &= \bar{Z}_{1i}(0) + \bar{M}_i(t) - \theta_{1i} \int_0^t \bar{Z}_{1i}(s)ds - \mu_{1i} \int_0^t \bar{Z}_{1i}(s)ds, \\
\bar{Z}_{2i}(t) &= \bar{Z}_{2i}(0) + \mu_{1i} \int_0^t \bar{Z}_{1i}(s)ds - \mu_2 \int_0^t \bar{Z}_{2i}(s)ds,
\end{align*}
where the matching process is given by
\begin{equation*}
\bar{M}(t) = \bar{M}_1(t) + \bar{M}_2(t) = \sup_{0 \le s \le t} \left\lceil\frac{\tilde{X}(s) + \tilde{Y}(s) - \sqrt{(\tilde{X}(s) - \tilde{Y}(s))^2 + 4(\bar{\mu}_{11}+\bar{\mu}_{12})}}{2}\right\rceil^+,
\end{equation*}
where $\tilde{X}_i(t) = \bar{Q}_i(0) + \lambda_i t - \theta_{0i} \int_0^t \bar{Q}_i(s)ds$ for $i=1,2$, $\tilde{X}(t) = \tilde{X}_1(t) + \tilde{X}_2(t)$, and
$\tilde{Y}(t) = \bar{Z}_0(t) + \sum_{i=1}^{2}(\theta_{1i} \int_0^t \bar{Z}_{1i}(s)ds + \mu_2 \int_0^t \bar{Z}_{2i}(s)ds)$.
\end{theorem}

\begin{proof}{Proof of Theorem~\ref{thm:EC1}}
For each $n$, we define the residual terms as follows:
\begin{align*}
B^n_{1i}(t) &= (\bar{A}^n_i(t) - \lambda_i t) - (\bar{R}^n_{0i}(t) - \theta_{0i} \int_0^t \bar{Q}^n_i(s)ds) - (\bar{M}^n_i(t) - \tilde{M}^n_i(t)), \\
B^n_2(t) &= \sum_{i=1}^{2} \{(\bar{R}^n_{1i}(t) - \theta_{1i} \int_0^t \bar{Z}^n_{1i}(s)ds) + (\bar{D}^n_{2i}(t) - \mu_2 \int_0^t \bar{Z}^n_{2i}(s)ds)\} \\
&\quad - \sum_{i=1}^{2}(\bar{M}^n_i(t) - \tilde{M}^n_i(t)), \\
B^n_{3i}(t) &= -(\bar{R}^n_{1i}(t) - \theta_{1i} \int_0^t \bar{Z}^n_{1i}(s)ds) - (\bar{D}^n_{1i}(t) - \mu_{1i} \int_0^t \bar{Z}^n_{1i}(s)ds) \\
&\quad + (\bar{M}^n_i(t) - \tilde{M}^n_i(t)), \\
B^n_{4i}(t) &= (\bar{D}^n_{1i}(t) - \mu_{1i} \int_0^t \bar{Z}^n_{1i}(s)ds) - (\bar{D}^n_{2i}(t) - \mu_2 \int_0^t \bar{Z}^n_{2i}(s)ds).
\end{align*}

With these definitions, the scaled stochastic processes $(\bar{Q}^n_i, \bar{Z}^n_0, \bar{Z}^n_{1i}, \bar{Z}^n_{2i}, i=1,2)$ and residual terms $(B^n_{1i}, B^n_2, B^n_{3i}, B^n_{4i}, i=1,2)$ jointly satisfy the system of equations \eqref{eq:EC39}--\eqref{eq:EC42}.

By the functional strong law of large numbers, each centered arrival and service process converges to zero in probability:
\begin{align*}
\lim_{n \to \infty} \mathbb{P}\left(\sup_{t \le T} |\bar{A}^n_i(t) - \lambda_i t| > \epsilon\right) &= 0, \quad i=1,2, \\
\lim_{n \to \infty} \mathbb{P}\left(\sup_{t \le T} \left|\bar{R}^n_{0i}(t) - \theta_{0i} \int_0^t \bar{Q}^n_i(s)ds\right| > \epsilon\right) &= 0, \quad i=1,2, \\
\lim_{n \to \infty} \mathbb{P}\left(\sup_{t \le T} \left|\bar{R}^n_{1i}(t) - \theta_{1i} \int_0^t \bar{Z}^n_{1i}(s)ds\right| > \epsilon\right) &= 0, \quad i=1,2, \\
\lim_{n \to \infty} \mathbb{P}\left(\sup_{t \le T} \left|\bar{D}^n_{2i}(t) - \mu_2 \int_0^t \bar{Z}^n_{2i}(s)ds\right| > \epsilon\right) &= 0, \quad i=1,2.
\end{align*}

From Lemma~\ref{lem:EC2}, we have
\begin{align*}
\lim_{n \to \infty} \mathbb{P}\left(\sup_{t \le T} |\bar{M}^n(t) - \tilde{M}^n(t)| > \epsilon\right) &= 0, \\
\lim_{n \to \infty} \mathbb{P}\left(\sup_{t \le T} \left|\bar{D}^n_{1i}(t) - \mu_{1i} \int_0^t \bar{Z}^n_{1i}(s)ds\right| > \epsilon\right) &= 0, \quad i=1,2.
\end{align*}

Therefore, for any $\epsilon > 0$,
\begin{equation*}
\lim_{n \to \infty} \mathbb{P}\left(\sup_{t \le T} \|(B^n_{1i}, B^n_2, B^n_{3i}, B^n_{4i}, i=1,2)\|_T > \epsilon\right) = 0,
\end{equation*}
which implies that $(B^n_{1i}, B^n_2, B^n_{3i}, B^n_{4i}, i=1,2)$ converges to $(0,0,0,0,0,0,0)$ in probability.

Since the mapping $f$ defined by equations \eqref{eq:EC39}--\eqref{eq:EC42} is Lipschitz continuous (Proposition~\ref{prop:EC1}), and $(B^n_{1i}, B^n_2, B^n_{3i}, B^n_{4i}, i=1,2) \to (0,0,0,0,0,0,0)$ in probability, we conclude that
\begin{equation*}
(\bar{Q}^n_i, \bar{Z}^n_0, \bar{Z}^n_{1i}, \bar{Z}^n_{2i}, i=1,2) \to (\bar{Q}_i, \bar{Z}_0, \bar{Z}_{1i}, \bar{Z}_{2i}, i=1,2)
\end{equation*}
in probability, where the limit satisfies the fluid model equations with the matching process.
\Halmos
\end{proof}

\section{Fluid Model: Existence, Uniqueness, and Equilibrium Convergence}
\label{sec:EC of fluid model convergence}

\begin{align}
q_i(t) &= q_i(0) + \lambda_i t - \theta_{0i}\int_0^t q_i(s)ds - m_i(t), \quad i=1,2, \label{eq:EC.53}  \\
z_0(t) &= z_0(0) + \sum_{j=1}^{2}\int_0^t (\theta_{1j}z_{1j}(s) + \mu_2z_{2j}(s))ds - \sum_{j=1}^{2} m_j(t), \label{eq:EC.54} \\
z_{1i}(t) &= z_{1i}(0) + m_i(t) - (\theta_{1i} + \mu_{1i})\int_0^t z_{1i}(s)ds, \quad i=1,2, \label{eq:EC.55}  \\
z_{2i}(t) &= z_{2i}(0) + \mu_{1i}\int_0^t z_{1i}(s)ds - \mu_2\int_0^t z_{2i}(s)ds, \quad i=1,2. \label{eq:EC.56} 
\end{align}

Throughout this section, we set the spatial matching parameters as $\alpha_1 = \alpha_2 = \alpha$.
This specific condition is adopted to ensure consistency with the derivation of matching rates under threshold policies previously established in Section~\ref{sec:ec_stochastic_to_fluid} (see Lemma~\ref{lem:EC1} and Theorem~\ref{thm:EC4}).
The fluid model for the two-type customer system, under this assumption, is defined by equations (EC.53)-(EC.56).

\begin{proposition} \label{prop:EC3}
The fluid model for the two-type customer system, as defined by equations (EC.53)-(EC.56) with the threshold matching policy where matching for type $i \in \{1,2\}$ occurs if its potential matching intensity $\mu_{1i}(t) = C q_i(t)^{\alpha} z_0(t)^{\alpha}$ meets or exceeds a threshold $\mu_{1i}^*$, has a unique solution trajectory $(q_i(t), z_0(t), z_{1i}(t), z_{2i}(t))$ for $t \ge 0$ given initial conditions.
\end{proposition}

\begin{proof}{Proof of Proposition~\ref{prop:EC3}}
The proof analyzes the system dynamics in stages, based on whether the matching thresholds are met.

\textbf{Stage 1: No matching thresholds are met.}
Assume the system starts in or enters a state where $\mu_{1i}(t) < \mu_{1i}^*$ for both $i=1,2$.
In this stage, no matching occurs, so the matching rates $m_1(t) = 0$ and $m_2(t) = 0$.
Assume $\mu_{1i}(t) = \mu_{1i}, i=1,2$ remains below the threshold, and the system dynamics \eqref{eq:fluid_q1}-\eqref{eq:policy_cond2} simplify to:
\begin{align}
q_i(t) &= q_i(0) + \lambda_i t - \theta_{0i} \int_0^t q_i(s)ds, \quad i=1,2, \label{eq:EC.57} \\
z_0(t) &= z_0(0) + \int_0^t \sum_{j=1}^{2}(\theta_{1j}z_{1j}(s) + \mu_2z_{2j}(s))ds, \label{eq:EC.58} \\
z_{1i}(t) &= z_{1i}(0) e^{-(\theta_{1i} + \mu_{1i})t}, \quad i=1,2, \label{eq:EC.59} \\
z_{2i}(t) &= z_{2i}(0) + \frac{\mu_{1i}}{\theta_{1i} + \mu_{1i}} z_{1i}(0) (1 - e^{-(\theta_{1i} + \mu_{1i})t}) - \mu_2\int_0^t z_{2i}(s)ds, \quad i=1,2. \label{eq:EC.60}
\end{align}
The derivatives of the right-hand sides of these equations are linear in the state variables, hence Lipschitz continuous.
By the Picard-Lindel\"of theorem~\citep{coddington1955theory}, a unique solution exists for this stage.
Typically, $q_i(t)$ evolves towards $\lambda_i/\theta_{0i}$ (if $\theta_{0i}>0$),
$z_{1i}(t) \to 0$, $z_{2i}(t) \to 0$, and $z_0(t) \to 1$
(due to driver conservation, assuming $\sum (z_{0} + z_{1j} + z_{2j}) = 1$).
Consequently, $\mu_{1i}(t)$ tends towards $C (\lambda_i/\theta_{0i})^{\alpha}$.
If system parameters are such that $C (\lambda_i/\theta_{0i})^{\alpha} > \mu_{1i}^*$
for some $i$, then $\mu_{1i}(t)$ for that type $i$ will eventually reach its threshold
$\mu_{1i}^*$. Let $t_1$ be the first time this occurs.

\textbf{Stage 2: Both matching thresholds are met and active.}
We set the spatial matching parameters as $\alpha_1 = \alpha_2 = \alpha$ for simplicity, and the case that $\alpha_1 \neq \alpha_2$ can be proven similarly.
Consider the scenario where both thresholds are met and active for $t \ge t_1$, i.e., $\mu_{1i}(t) = C q_i(t)^{\alpha} z_0(t)^{\alpha} = \mu_{1i}^*$ for $i=1,2$, which implies $\frac{d}{dt} (C q_i(t)^{\alpha} z_0(t)^{\alpha}) = 0$.
The condition $\frac{d}{dt} (C q_i(t)^{\alpha} z_0(t)^{\alpha}) = 0$ for $i=1,2$ implies that the logarithmic derivative is zero. Since $C$ and $\alpha$ are constants, and assuming $q_i(t) > 0$ and $z_0(t) > 0$:
\begin{equation*}
  \frac{d}{dt} (\alpha \ln q_i(t) + \alpha \ln z_0(t)) = 0
\end{equation*}
\begin{equation*}
  \alpha \left( \frac{1}{q_i(t)}\frac{dq_i}{dt} + \frac{1}{z_0(t)}\frac{dz_0}{dt} \right) = 0
\end{equation*}
Multiplying by $q_i(t)z_0(t)/\alpha$ (assuming $\alpha \neq 0$) gives
\begin{equation*}
  z_0(t)\frac{dq_i}{dt} + q_i(t)\frac{dz_0}{dt} = 0, \quad i=1,2.
\end{equation*}
Now, we substitute the expressions for the derivatives $\frac{dq_i}{dt}$ and $\frac{dz_0}{dt}$ from equations (EC.53) and (EC.54), note that $m_i(t)$ is monotonically increasing and non-negative so $\pi_i(t) \triangleq m_i'(t)$ exists a.e., $i=1,2$:
\begin{align*} 
    \frac{dq_i}{dt} &= \lambda_i - \theta_{0i}q_i(t) - \pi_i(t), \\ 
    \frac{dz_0}{dt} &= \sum_{j=1}^{2}(\theta_{1j}z_{1j}(t) + \mu_2z_{2j}(t)) - \sum_{j=1}^{2}\pi_j(t).
\end{align*}
Let $s(t) = \sum_{j=1}^{2}(\theta_{1j}z_{1j}(t) + \mu_2z_{2j}(t))$. Then $\frac{dz_0}{dt} = s(t) - (\pi_1(t) + \pi_2(t))$.

Consider the case for $i=1$:
\begin{equation*}
  z_0(t)\left( \lambda_1 - \theta_{01}q_1(t) - \pi_1(t) \right) + q_1(t)\left( s(t) - \pi_1(t) - \pi_2(t) \right) = 0.
\end{equation*}
Expanding,
\begin{equation*}
  z_0(t)\lambda_1 - z_0(t)\theta_{01}q_1(t) - z_0(t)\pi_1(t) + q_1(t)s(t) - q_1(t)\pi_1(t) - q_1(t)\pi_2(t) = 0.
\end{equation*}
Collecting terms involving $\pi_1(t)$ and $\pi_2(t)$ on the left-hand side,
\begin{equation*}
  z_0(t)\pi_1(t) + q_1(t)\pi_1(t) + q_1(t)\pi_2(t) = z_0(t)\lambda_1 - z_0(t)\theta_{01}q_1(t) + q_1(t)s(t).
\end{equation*}
Factoring yields
\begin{equation*}
  (z_0(t) + q_1(t))\pi_1(t) + q_1(t)\pi_2(t) = (\lambda_1 - \theta_{01}q_1(t))z_0(t) + s(t)q_1(t).
\end{equation*}
This is the first equation of the linear system.
Repeating the same calculation for $i=2$ produces the second equation, giving the following linear system for $\pi_1(t)$ and $\pi_2(t)$:
\begin{align*}
(z_0(t) + q_1(t))\pi_1(t) + q_1(t)\pi_2(t) &= (\lambda_1 - \theta_{01}q_1(t))z_0(t) + s(t)q_1(t), \\
q_2(t)\pi_1(t) + (z_0(t) + q_2(t))\pi_2(t) &= (\lambda_2 - \theta_{02}q_2(t))z_0(t) + s(t)q_2(t).
\end{align*}
The determinant of the coefficient matrix is $z_0(t)(z_0(t) + q_1(t) + q_2(t))$.
Since $\mu_{1i}(t) = \mu_{1i}^* > 0$ implies $q_i(t) > 0$ and $z_0(t) > 0$, the determinant is strictly positive.
Thus, we can uniquely solve for $\pi_1(t)$ and $\pi_2(t)$ as functions of the state variables, denoted $m_i^*(q_1, q_2, z_0, z_{11}, z_{12}, z_{21}, z_{22})$.
Substituting these $\pi_i^*(\cdot)$ into (EC.53)-(EC.56) yields a system of ODEs whose right-hand sides are continuous and, on any compact subset of the state space where $q_i, z_0$ are bounded away from zero, Lipschitz continuous.
The Picard-Lindelöf theorem then guarantees a unique solution for this stage.

If the system state evolves such that a matching condition $\mu_{1k}(t) < \mu_{1k}^*$ for some $k$, then $\pi_k(t)$ becomes 0, and the system dynamics adjust accordingly, potentially re-entering Stage 1 or a mixed stage.

\textbf{Stage 3: Mixed stage - one threshold active, one not.}
Without loss of generality, assume $\mu_{11}(t) = C q_1(t)^{\alpha} z_0(t)^{\alpha} = \mu_{11}^*$ (active) and $\mu_{12}(t) < \mu_{12}^*$ (inactive).
In this case, $\pi_2(t) = 0$. The condition for the active threshold $\mu_{11}^*$ is
\begin{equation*}
  z_0(t)\frac{dq_1}{dt} + q_1(t)\frac{dz_0}{dt} = 0.
\end{equation*}
Substituting the derivatives,
\begin{equation*}
  \frac{dq_1}{dt} = \lambda_1 - \theta_{01}q_1(t) - \pi_1(t),
\end{equation*}
\begin{equation*}
  \frac{dz_0}{dt} = \sum_{j=1}^{2}(\theta_{1j}z_{1j}(t) + \mu_2z_{2j}(t)) - \pi_1(t) \quad (\text{since } \pi_2(t)=0).
\end{equation*}
Let $s(t) = \sum_{j=1}^{2}(\theta_{1j}z_{1j}(t) + \mu_2z_{2j}(t))$.
The condition becomes
\begin{equation*}
  z_0(t)(\lambda_1 - \theta_{01}q_1(t) - \pi_1(t)) + q_1(t)(s(t) - \pi_1(t)) = 0,
\end{equation*}
\begin{equation*}
  z_0(t)\lambda_1 - z_0(t)\theta_{01}q_1(t) - z_0(t)\pi_1(t) + q_1(t)s(t) - q_1(t)\pi_1(t) = 0.
\end{equation*}
Solving for $\pi_1(t)$,
\begin{equation*}
  (z_0(t) + q_1(t))\pi_1(t) = z_0(t)\lambda_1 - z_0(t)\theta_{01}q_1(t) + q_1(t)s(t),
\end{equation*}
\begin{equation*}
  \pi_1(t) = \frac{(\lambda_1 - \theta_{01}q_1(t))z_0(t) + s(t)q_1(t)}{z_0(t) + q_1(t)}.
\end{equation*}
Since $\mu_{11}^* > 0$, $q_1(t) > 0$ and $z_0(t) > 0$, so the denominator $z_0(t) + q_1(t) > 0$.
Thus, $\pi_1(t)$ is uniquely determined as a function of the state variables.
Substituting $\pi_1(t)$ and $\pi_2(t)=0$ into (EC.53)-(EC.56) yields a system of ODEs.
Similar to Stage 2, the right-hand sides are continuous and Lipschitz on relevant compact sets, guaranteeing a unique solution.
The system remains in this stage as long as $\mu_{11}(t) = \mu_{11}^*$ and $\mu_{12}(t) < \mu_{12}^*$.
Transitions occur if $\mu_{12}(t)$ reaches $\mu_{12}^*$ (to Stage 2) or if the conditions for maintaining $\mu_{11}(t) = \mu_{11}^*$ (e.g., $\pi_1(t) \ge 0$) are violated, potentially leading to Stage 1.

\textbf{Overall Existence and Uniqueness.}
The unique solutions within each stage, combined with deterministic conditions for transitioning between stages (i.e., $\mu_{1i}(t)$ reaching or falling below $\mu_{1i}^*$), establish the existence and uniqueness of the overall fluid trajectory.
The matching rates $\pi_i(t)$ are defined piecewise based on these conditions, ensuring a well-defined path for the system state.
\Halmos
\end{proof}

\begin{theorem} \label{thm:EC4}
Consider the fluid model where matching for type $i \in \{1,2\}$ occurs if its potential matching intensity $\mu_{1i}(t) = C q_i(t)^{\alpha} z_0(t)^{\alpha}$ meets or exceeds a positive threshold $\mu_{1i}^* > 0$.
If system parameters allow for a physical equilibrium where matching is active for both types (meaning the derived matching rates $\pi_i^*$ are non-negative and $z_0^* > 0$), then there exists a unique non-negative solution $(q_1^*, q_2^*, z_0^*, z_{11}^*, z_{12}^*, z_{21}^*, z_{22}^*)$ with associated non-negative matching rates $(\pi_1^*, \pi_2^*)$ to the equilibrium equations:
\begin{align}
\lambda_i &= \theta_{0i}q_i^* + \theta_{1i}z_{1i}^* + \mu_2z_{2i}^*, && i=1,2, \label{eq:thm_eq_1} \\
\mu_{1i}^*z_{1i}^* &= \mu_2z_{2i}^*, && i=1,2, \label{eq:thm_eq_2} \\
1 &= z_0^* + \sum_{j=1}^{2}(z_{1j}^* + z_{2j}^*), && \label{eq:thm_eq_3} \\
\mu_{1i}^* &= C (q_i^*)^{\alpha} (z_0^*)^{\alpha} , && i=1,2. \label{eq:thm_eq_4}
\end{align}
This solution is the unique physical equilibrium point of the fluid model when both matching thresholds $\mu_{1i}^*$ are active.
Furthermore, the fluid process $(q_i(t), z_0(t), z_{1i}(t), z_{2i}(t))$ for $i=1,2$, converges exponentially fast to this equilibrium point as $t \to \infty$.
\end{theorem}

\begin{proof}{Proof of Theorem~\ref{thm:EC4}}
This proof addresses the existence, uniqueness, and convergence to an equilibrium point $(q_1^*, q_2^*, z_0^*, z_{11}^*, z_{12}^*, z_{21}^*, z_{22}^*)$ where both matching thresholds are active.
This means the matching intensity for type $i$, $\mu_{1i}(t) = C q_i(t)^{\alpha} z_0(t)^{\alpha}$, is equal to a predefined positive threshold $\mu_{1i}^*$ for $i=1,2$.

\textbf{1. Existence and Uniqueness of a Solution to Equilibrium Equations}

At an equilibrium point, all time derivatives in the fluid model equations (EC.53)-(EC.56) must be zero. This yields the following system of algebraic equations:
\begin{align}
\lambda_i &= \theta_{0i}q_i^* + \pi_i^*, \quad i=1,2, \label{eq:ThmEC4.Eq1}  \\
\sum_{j=1}^{2}\pi_j^* &= \sum_{j=1}^{2}(\theta_{1j}z_{1j}^* + \mu_{2}z_{2j}^*), \label{eq:ThmEC4.Eq2} \\
\pi_i^* &= (\theta_{1i} + \mu_{1i}^*)z_{1i}^*, \quad i=1,2, \label{eq:ThmEC4.Eq3}  \\
\mu_{1i}^*z_{1i}^* &= \mu_{2}z_{2i}^*, \quad i=1,2. \label{eq:ThmEC4.Eq4} 
\end{align}
Additionally, the conservation of drivers (total proportion is 1) requires:
\begin{equation}
1 = z_0^* + \sum_{j=1}^{2}(z_{1j}^* + z_{2j}^*) \label{eq:ThmEC4.Eq5} 
\end{equation}
For an equilibrium with active matching thresholds as assumed:
\begin{equation}
C (q_i^*)^{\alpha} (z_0^*)^{\alpha} = \mu_{1i}^*, \quad i=1,2 \label{eq:ThmEC4.Eq6} 
\end{equation}
where $\mu_{1i}^* > 0$ are the predefined thresholds for the matching intensities and we assume $\alpha_1 = \alpha_2 = \alpha$ for simplicity. 
Equation \eqref{eq:ThmEC4.Eq6} implies $q_i^* > 0$ and $z_0^* > 0$.

From \eqref{eq:ThmEC4.Eq6}, we express $q_i^*$ as $q_i^* = (\mu_{1i}^*/C)^{1/\alpha} (z_0^*)^{-1}$.
From \eqref{eq:ThmEC4.Eq3} and \eqref{eq:ThmEC4.Eq4}, assuming $\theta_{1j} + \mu_{1j}^* > 0$ and $\mu_2 > 0$, we can express the assigned drivers and busy drivers as: $z_{1j}^* = \pi_j^*/(\theta_{1j} + \mu_{1j}^*)$ and $z_{2j}^* = (\mu_{1j}^*/\mu_2)z_{1j}^* = \frac{\mu_{1j}^*}{\mu_2} \frac{\pi_j^*}{\theta_{1j} + \mu_{1j}^*}$.
Substituting these expressions for $z_{1j}^*$ and $z_{2j}^*$ into the conservation equation \eqref{eq:ThmEC4.Eq5} gives
$1 = z_0^* + \sum_{j=1}^{2} \pi_j^* \left( \frac{1}{\theta_{1j} + \mu_{1j}^*} (1 + \frac{\mu_{1j}^*}{\mu_2}) \right)$.
Let $k_j = \frac{1 + \mu_{1j}^*/\mu_2}{\theta_{1j} + \mu_{1j}^*}$. Then $1 = z_0^* + \sum_{j=1}^{2} k_j \pi_j^*$.
Using $\pi_j^* = \lambda_j - \theta_{0j}q_j^*$ from \eqref{eq:ThmEC4.Eq1}, and then substituting the expression for $q_j^*$:
$1 = z_0^* + \sum_{j=1}^{2} k_j \left( \lambda_j - \theta_{0j}(\mu_{1j}^*/C)^{1/\alpha} (z_0^*)^{-1} \right)$.
Multiplying by $z_0^*$ (since $z_0^* > 0$ at this type of equilibrium):
$z_0^* = (z_0^*)^2 + \sum_{j=1}^{2} k_j \lambda_j z_0^* - \sum_{j=1}^{2} k_j \theta_{0j}(\mu_{1j}^*/C)^{1/\alpha}$.
Rearranging gives a quadratic equation for $z_0^*$:
\begin{equation}
(z_0^*)^2 + \left( \sum_{j=1}^{2} k_j \lambda_j - 1 \right) z_0^* - \sum_{j=1}^{2} k_j \theta_{0j}(\mu_{1j}^*/C)^{1/\alpha} = 0. \label{eq:ThmEC4.Eq7}
\end{equation}
Let $A=1$, $B = \sum_{j=1}^{2} k_j \lambda_j - 1$, and $D = - \sum_{j=1}^{2} k_j \theta_{0j}(\mu_{1j}^*/C)^{1/\alpha}$.
Assume all system parameters ($\lambda_j, \theta_{0j}, \theta_{1j}, \mu_{1j}, \mu_2, C, \mu_{1j}^*$) are positive.
Then $k_j > 0$.
If $\theta_{0j} > 0$ and $\mu_{1j}^* > 0$ for at least one $j$, then $D < 0$.
If $D < 0$, the quadratic equation $A(z_0^*)^2 + B z_0^* + D = 0$ has exactly one positive real root given by $z_0^* = (-B + \sqrt{B^2 - 4AD})/(2A)$.
This provides a unique candidate for $z_0^*$.
(If $D=0$, which would imply $\theta_{0j}(\mu_{1j}^*/C)^{1/\alpha}=0$ for all $j$, then $z_0^*(z_0^*+B)=0$.
A positive solution $z_0^* = -B$ exists if $B<0$.)

Given this unique positive $z_0^*$:
\begin{itemize}
    \item $q_i^* = (\mu_{1i}^*/C)^{1/\alpha} (z_0^*)^{-1}$ is uniquely determined and positive.
    \item $\pi_i^* = \lambda_i - \theta_{0i}q_i^*$ is uniquely determined. For a physical equilibrium of this type (both thresholds active), we require $\pi_i^* \ge 0$. This imposes conditions on the system parameters, namely $\lambda_i \ge \theta_{0i}q_i^*$.
    \item $z_{1i}^* = \pi_i^*/(\theta_{1i} + \mu_{1i}^*)$ is uniquely determined and non-negative if $\pi_i^* \ge 0$.
    \item $z_{2i}^* = (\mu_{1i}^*/\mu_2)z_{1i}^*$ is uniquely determined and non-negative if $z_{1i}^* \ge 0$.
\end{itemize}
Thus, if system parameters ensure $D<0$ (or $D=0$ and $B<0$ for a positive $z_0^*$) and the resulting matching rates $m_i^*$ are non-negative, there exists a unique non-negative solution $(q_1^*, q_2^*, z_0^*, z_{11}^*, z_{12}^*, z_{21}^*, z_{22}^*)$ to the equilibrium equations \eqref{eq:ThmEC4.Eq1}-\eqref{eq:ThmEC4.Eq6}.

\textbf{2. The Solution is the Unique Equilibrium of the Fluid Model (under active threshold conditions)}

The solution derived in Part 1, by its construction, satisfies the conditions for an equilibrium point of the fluid model where all time derivatives are zero and both matching thresholds $\mu_{1i}^*$ are active.
Proposition~\ref{prop:EC3} establishes that the fluid model equations (EC.53)-(EC.56), under the specified threshold matching policy, have a unique solution trajectory given any initial conditions.
As the system parameters and initial conditions lead the system to an equilibrium where both matching thresholds are active (i.e., the system operates in Stage 2, as described in the proof of Proposition~\ref{prop:EC3}, for $t \to \infty$), then this equilibrium point must be the unique point found in Part 1.

\textbf{3. Convergence to the Equilibrium}

To show convergence of the fluid process $(q_i(t), z_0(t), z_{1i}(t), z_{2i}(t))$ to the equilibrium point $X^* = (q_1^*, q_2^*, z_0^*, z_{11}^*, z_{12}^*, z_{21}^*, z_{22}^*)$ found in Part 1, we analyze the dynamics of the system under the assumption that both matching thresholds are active.
This means $\mu_{1i}(t) = C q_i(t)^{\alpha} z_0(t)^{\alpha} = \mu_{1i}^* > 0$ for $i=1,2$, so $q_i(t)z_0(t) = \sqrt[\alpha]{\mu_{1i}^*/C} =: K_i$, where $K_i$ is a positive constant for $i=1,2$.
This implies $q_i'(t)z_0(t) + q_i(t)z_0'(t) = 0$.

The argument for convergence proceeds in two main steps:
(a) First, we argue that $z_0(t)$ (and consequently $q_i(t)=K_i/z_0(t)$) converges to its equilibrium value $z_0^*$ (and $q_i^*$).
(b) Second, given the convergence of $z_0(t)$ and $q_i(t)$ to their equilibrium values, we analyze the subsystem for $(z_{11}(t), z_{12}(t), z_{21}(t), z_{22}(t))$. We show that this subsystem, when $z_0, q_i$ are at equilibrium, is governed by a system of linear autonomous ODEs whose Jacobian matrix (denoted $J$) has all eigenvalues with strictly negative real parts. This implies exponential convergence for these components, which is key to the "exponentially fast" claim for the overall system convergence.

\textbf{(a) Convergence of $z_0(t)$ and $q_i(t)$:}
Let $A_{tot}(t) = \sum_{j=1}^{2} (\lambda_j - \theta_{0j}q_j(t))$ and 
$B_{state}(t) = \sum_{j=1}^{2}(\theta_{1j}z_{1j}(t) + \mu_2z_{2j}(t))$. 
The total matching rate $\Pi(t) = \pi_1(t)+\pi_2(t)$, as derived in Proposition~\ref{prop:EC3} (Stage 2), 
is a weighted average of $A_{tot}(t)$ and $B_{state}(t)$ when both thresholds are active:
\begin{equation*}
  \Pi(t)(z_0(t) + \sum_{j=1}^{2} q_j(t)) = A_{tot}(t)z_0(t) + B_{state}(t)(\sum_{j=1}^{2} q_j(t))
\end{equation*}
This implies $\Pi(t)$ lies between $A_{tot}(t)$ and $B_{state}(t)$.
The dynamics for $z_0(t)$ are $z_0'(t) = B_{state}(t) - \Pi(t)$.
Let $Q(t) = \sum_{j=1}^{2} q_j(t)$. From $q_j'(t)z_0(t) + q_j(t)z_0'(t) = 0$, 
summing over $j$ gives $Q'(t)z_0(t) + Q(t)z_0'(t) = 0$.
Also, $Q'(t) = \sum q_j'(t) = \sum (\lambda_j - \theta_{0j}q_j(t) - \pi_j(t)) = A_{tot}(t) - \Pi(t)$.
Thus, $z_0'(t)$ and $Q'(t)$ have opposite signs if $A_{tot}(t) \neq B_{state}(t)$ (assuming $z_0(t)>0, Q(t)>0$).

   \textit{Case i: $A_{tot}(t) - B_{state}(t)$ has a fixed sign for $t \ge t_{start}$ (e.g., $A_{tot}(t) \ge B_{state}(t)$).}
       Then $A_{tot}(t) \ge \Pi(t) \ge B_{state}(t)$. This implies $z_0'(t) = B_{state}(t) - \Pi(t) \le 0$. 
       Since $z_0(t)$ is bounded below by $z_0^* > 0$ (as per Part 1, assuming a physical equilibrium 
       with $z_0^*>0$), $z_0(t)$ is non-increasing and converges to some limit $\tilde{z}_0 \ge z_0^*$. 
       Consequently, $q_i(t) = K_i/z_0(t)$ is non-decreasing and converges to $\tilde{q}_i = K_i/\tilde{z}_0$.
       The state variables $z_{1j}(t), z_{2j}(t)$ are bounded (e.g., between 0 and 1). As $t \to \infty$, 
       if the system converges to an equilibrium, then $z_0'(t) \to 0$, 
       which implies $B_{state}(t) \to \Pi(t)$. Also, $A_{tot}(t) \to \Pi(t)$. 
       This means $A_{tot}(\tilde{q}) = B_{state}(\tilde{z}_{1},\tilde{z}_{2})$. 
       Given the uniqueness of the overall equilibrium $X^*$ (from Part 1) where all derivatives are zero, it must be that this limit point is indeed $X^*$, so $\tilde{z}_0 = z_0^*$ and $\tilde{q}_i = q_i^*$.
       Thus, $z_0(t) \to z_0^*$ and $q_i(t) \to q_i^*$.
       The case $A_{tot}(t) \le B_{state}(t)$ is analogous, leading to $z_0(t)$ being non-decreasing and converging to $z_0^*$.

 \textit{Case ii: $A_{tot}(t) - B_{state}(t)$ changes sign or reaches 0.}
    Suppose at some time $t_1$, $f(t_1) = A_{tot}(t_1) - B_{state}(t_1) = 0$.
    If $X(t_1)$ is the equilibrium point $X^*$, then $f(t) \equiv 0$ for $t \ge t_1$ as all derivatives 
    are zero.
    If $X(t_1) \neq X^*$, then $q_j'(t_1)=0$ and $z_0'(t_1)=0$ (as shown in the reasoning for Case i 
    when $A_{tot}=B_{state}=\Pi$). However, since $X(t_1) \neq X^*$, at least one of $z_{1j}'(t_1)$ 
    or $z_{2j}'(t_1)$ must be non-zero. Let $Z'_{sub}(t_1)$ be the vector of these derivatives.
    We analyze the behavior of $f(t)$ in the neighborhood of $t_1$ by looking at its derivatives.
    $f'(t) = A_{tot}'(t) - B_{state}'(t)$. At $t=t_1$:
    $A_{tot}'(t_1) = \sum_{j=1}^{2} (-\theta_{0j}q_j'(t_1)) = 0$.
    So, $f'(t_1) = -B_{state}'(t_1) = -\sum_{j=1}^{2}(\theta_{1j}z_{1j}'(t_1) + \mu_2z_{2j}'(t_1))$.

    If $f'(t_1) \neq 0$: Since $f(t_1)=0$, $f(t)$ will change sign locally around $t_1$.
    Specifically, for $t > t_1$ and $t$ close to $t_1$, $f(t)$ will have the same sign as $f'(t_1)$.
    For example, if $f'(t_1) > 0$, then $f(t) > 0$ for $t \in (t_1, t_1+\delta)$ for some $\delta > 0$.
    This means $A_{tot}(t) > B_{state}(t)$ in this interval, which is the scenario of Case i.
    Consequently, $z_0(t)$ becomes monotonic (non-increasing in this example) and converges to $z_0^*$.

    If $f'(t_1) = 0$: We then examine $f''(t_1)$.
    $A_{tot}''(t_1) = \sum_{j=1}^{2} (-\theta_{0j}q_j''(t_1))$.
    The condition $q_j'(t)z_0(t) + q_j(t)z_0'(t) = 0$, when differentiated and evaluated at $t_1$ (where $q_j'(t_1)=0, z_0'(t_1)=0$), yields $q_j''(t_1)z_0(t_1) + q_j(t_1)z_0''(t_1) = 0$.
    Also, $z_0'(t) = B_{state}(t) - \Pi(t)$. When $A_{tot}(t_1)=B_{state}(t_1)$, $\Pi(t_1)=A_{tot}(t_1)$,
    so $z_0'(t_1) = B_{state}(t_1) - A_{tot}(t_1) = -f(t_1)$.
    Thus $z_0''(t_1) = -f'(t_1)$. Since we are in the subcase $f'(t_1)=0$, we have $z_0''(t_1)=0$.
    If $z_0(t_1) > 0$, then $q_j''(t_1)z_0(t_1)=0$ which implies $q_j''(t_1)=0$. So $A_{tot}''(t_1)=0$.
    Therefore, $f''(t_1) = -B_{state}''(t_1) = -\sum_{j=1}^{2}(\theta_{1j}z_{1j}''(t_1) + \mu_2z_{2j}''(t_1))$.
    If $f''(t_1) \neq 0$: Since $f(t_1)=0$ and $f'(t_1)=0$, $f(t)$ has a local extremum at $t_1$.
    For $t \neq t_1$ near $t_1$, $f(t)$ will have the sign of $f''(t_1)$.
    For example, if $f''(t_1) > 0$, then $f(t) > 0$ for $t \in (t_1-\delta, t_1) \cup (t_1, t_1+\delta)$.
    This means $A_{tot}(t) \ge B_{state}(t)$ locally, which again leads to the monotonic behavior of $z_0(t)$ as in Case i (allowing for $A_{tot}(t)=B_{state}(t)$ at $t_1$).

    This process can be continued.
    If $f(t_1)=f'(t_1)=\dots=f^{(k-1)}(t_1)=0$ and $f^{(k)}(t_1) \neq 0$ for some $k \ge 1$:
    Then $f(t) \approx f^{(k)}(t_1)\frac{(t-t_1)^k}{k!}$ for $t$ near $t_1$.
    If $k$ is odd, $f(t)$ changes sign at $t_1$.
    For $t>t_1$ (locally), $f(t)$ has a fixed sign (that of $f^{(k)}(t_1)$), which leads to Case i behavior.
    If $k$ is even, $f(t)$ does not change sign at $t_1$; it has the sign of $f^{(k)}(t_1)$ for $t \neq t_1$ locally, which also leads to Case i behavior (where $A_{tot}(t)$ is either $\ge B_{state}(t)$ or $\le B_{state}(t)$ locally).
    In both situations, for $t > t_1$ and $t$ sufficiently close to $t_1$, $f(t)$ maintains a fixed sign (or is zero only at $t_1$ if $k$ is even).
    This ensures $z_0(t)$ is monotonic in $(t_1, t_1+\delta)$ and thus converges to $z_0^*$.

    The only scenario not covered is if $f^{(k)}(t_1)=0$ for all $k \ge 1$.
    This would imply $f(t) \equiv 0$ in a neighborhood of $t_1$.
    If $A_{tot}(t) \equiv B_{state}(t)$ in a neighborhood, then $z_0'(t) = -f(t) \equiv 0$, so $z_0(t)$ is constant.
    This implies $q_j(t) = K_j/z_0(t)$ are also constant. So $q_j'(t) \equiv 0$.
    Thus $A_{tot}(t)$ is constant, which means $B_{state}(t)$ must also be constant.
    $B_{state}'(t) = \sum_{j=1}^{2}(\theta_{1j}z_{1j}'(t) + \mu_2z_{2j}'(t)) \equiv 0$.
    The dynamics for $Z_{sub}(t)=(z_{11}(t), z_{12}(t), z_{21}(t), z_{22}(t))$ are $Z'_{sub}(t) = J_{sub}(Z_{sub}(t) - Z_{sub,eq})$, where $J_{sub}$ and $Z_{sub,eq}$ are constant because $q_j, z_0$ are constant.
    If $Z_{sub}(t_1) \neq Z_{sub,eq}$ (which must be true if $X(t_1) \neq X^*$), then $Z_{sub}(t) - Z_{sub,eq}$ is a non-zero solution to $x' = J_{sub}x$.
    The condition $B_{state}'(t) \equiv 0$ means $v_B \cdot Z'_{sub}(t) \equiv 0$, where $v_B = (\theta_{11}, \theta_{12}, \mu_{2}, \mu_{2})$ (matching the order of components in $Z_{sub}$).
    So $v_B \cdot J_{sub}(Z_{sub}(t) - Z_{sub,eq}) \equiv 0$, which implies $v_B J_{sub}^k (Z_{sub}(t_1) - Z_{sub,eq}) = 0$ for all $k \ge 0$.
    For a stable and observable system $(J_{sub}, v_B)$, this implies $Z_{sub}(t_1) - Z_{sub,eq} = 0$ (observability means that $\cap_{k=0}^{3} \ker(v_B J_{sub}^k) = \{0\}$).
    Thus $Z_{sub}(t_1) = Z_{sub,eq}$, meaning $X(t_1)$ satisfies all equilibrium conditions, so $X(t_1)=X^*$.
    This contradicts the assumption that $X(t_1) \neq X^*$.
    Therefore, if $X(t_1) \neq X^*$, there must be some $k \ge 1$ such that $f^{(k)}(t_1) \neq 0$.
    As shown above, this leads to $z_0(t)$ becoming locally monotonic and thus continuing its convergence towards $z_0^*$.

    Combining these arguments, if the system state $X(t)$ reaches a point where $A_{tot}(t) = B_{state}(t)$ but $X(t) \neq X^*$, the system dynamics will necessarily drive $A_{tot}(t) - B_{state}(t)$ to be non-zero (or at least one-sided) in the subsequent interval, re-establishing the conditions for monotonic convergence of $z_0(t)$ (and $q_i(t)$) as described in Case i.
    Sustained oscillations of $A_{tot}(t) - B_{state}(t)$ around zero are not possible because $X^*$ is the unique equilibrium.

Based on these arguments, we conclude that $z_0(t) \to z_0^*$ and $q_i(t) \to q_i^*$ as $t \to \infty$.

\textbf{(b) Exponential Convergence of $z_{1i}(t)$ and $z_{2i}(t)$ given $z_0(t) \to z_0^*, q_i(t) \to q_i^*$:}
Once $z_0(t)$ and $q_i(t)$ are close to their equilibrium values $z_0^*$ and $q_i^*$, we analyze the dynamics for the subsystem $Z_{sub}(t) = (z_{11}(t), z_{12}(t), z_{21}(t), z_{22}(t))^T$.
The governing equations become approximately:
\begin{align*} 
    z_{1i}'(t) &= \pi_i(z_0^*, q_1^*, q_2^*, Z_{sub}(t)) - (\theta_{1i} + \mu_{1i}^*)z_{1i}(t), \\ 
    z_{2i}'(t) &= \mu_{1i}^*z_{1i}(t) - \mu_2z_{2i}(t). 
\end{align*}
Let $s(Z_{sub}) = \sum_{j=1}^{2}(\theta_{1j}z_{1j} + \mu_2z_{2j})$.
The matching rates $\pi_i$ can be expressed as linear functions of $s(Z_{sub})$ when $z_0, q_1, q_2$ are fixed at their equilibrium values:
\begin{align*} 
    \pi_1(Z_{sub}) &= \frac{(\lambda_1 - \theta_{01}q_1^*)(z_0^*+q_2^*) + (s(Z_{sub})-(\lambda_2 - \theta_{02}q_2^*))q_1^*}{z_0^* + q_1^* + q_2^*} =: c_{11} + c_{12}s(Z_{sub}), \\ 
    \pi_2(Z_{sub}) &= \frac{(\lambda_2 - \theta_{02}q_2^*)(z_0^*+q_1^*) + (s(Z_{sub})-(\lambda_1 - \theta_{01}q_1^*))q_2^*}{z_0^* + q_1^* + q_2^*} =: c_{21} + c_{22}s(Z_{sub}). 
\end{align*}
The coefficients $c_{ij}$ are constants determined by the equilibrium values $z_0^*, q_1^*, q_2^*$ and system parameters $\lambda_k, \theta_{0k}$. Specifically, $c_{12} = \frac{q_1^*}{z_0^*+q_1^*+q_2^*}$ and $c_{22} = \frac{q_2^*}{z_0^*+q_1^*+q_2^*}$. Given $q_i^* > 0$ and $z_0^* > 0$, these satisfy $0 < c_{12} < 1$, $0 < c_{22} < 1$. Crucially, $c_{12}+c_{22} = \frac{q_1^*+q_2^*}{z_0^*+q_1^*+q_2^*} < 1$.
The system for $Z_{sub}$ is $Z_{sub}' = J Z_{sub} + K$, where $K = [c_{11}, c_{21}, 0, 0]^T$ and $J$ is the $4 \times 4$ Jacobian matrix:
\begin{equation*}
  J = \begin{pmatrix}
  c_{12}\theta_{11} - \theta_{11} - \mu_{11}^{*} & c_{12}\theta_{12} & c_{12}\mu_{2} & c_{12}\mu_{2} \\
  c_{22}\theta_{11} & c_{22}\theta_{12} - \theta_{12} - \mu_{12}^{*} & c_{22}\mu_{2} & c_{22}\mu_{2} \\
  \mu_{11}^{*} & 0 & -\mu_{2} & 0 \\
  0 & \mu_{12}^{*} & 0 & -\mu_{2}
  \end{pmatrix}.
\end{equation*}
The equilibrium for this subsystem is $Z_{sub}^* = (z_{11}^*, z_{12}^*, z_{21}^*, z_{22}^*)$.
We now prove that all eigenvalues of $J$ have strictly negative real parts.

\textit{Proof that eigenvalues of $J$ have strictly negative real parts:}
The matrix $J$ is a Metzler matrix (all off-diagonal elements are non-negative, as $c_{ij} \ge 0$ and all $\theta, \mu$ parameters are positive).
For a Metzler matrix, stability (all eigenvalues $\lambda$ satisfying $\text{Re}(\lambda) < 0$) is equivalent to the existence of a strictly positive vector $v \gg 0$ (all components $v_k > 0$) such that $Jv \ll 0$ (all components of $Jv$ are strictly negative).
Let $v = [v_1, v_2, v_3, v_4]^T$. The components of $Jv$ are:
\begin{align}
(Jv)_1 &= [(c_{12}-1)\theta_{11} - \mu_{11}^{*}]v_1 + c_{12}\theta_{12}v_2 + c_{12}\mu_{2}v_3 + c_{12}\mu_{2}v_4 \label{eq:Jv1}, \\
(Jv)_2 &= c_{22}\theta_{11}v_1 + [(c_{22}-1)\theta_{12} - \mu_{12}^{*}]v_2 + c_{22}\mu_{2}v_3 + c_{22}\mu_{2}v_4 \label{eq:Jv2}, \\
(Jv)_3 &= \mu_{11}^{*}v_1 - \mu_{2}v_3 \label{eq:Jv3}, \\
(Jv)_4 &= \mu_{12}^{*}v_2 - \mu_{2}v_4 \label{eq:Jv4}.
\end{align}
We need to find $v_k > 0$ such that $(Jv)_k < 0$ for $k=1,2,3,4$.

From \eqref{eq:Jv3}, we require $\mu_{11}^{*}v_1 - \mu_{2}v_3 < 0$. Choose $v_3 = (1+\epsilon)\frac{\mu_{11}^{*}}{\mu_{2}}v_1$ for some small $\epsilon > 0$. Then,
\begin{equation*}
  (Jv)_3 = \mu_{11}^{*}v_1 - \mu_{2}(1+\epsilon)\frac{\mu_{11}^{*}}{\mu_{2}}v_1 = \mu_{11}^{*}v_1 - (1+\epsilon)\mu_{11}^{*}v_1 = -\epsilon\mu_{11}^{*}v_1.
\end{equation*}
Since $v_1 > 0$ and $\epsilon > 0$, $(Jv)_3 < 0$.

Similarly, from \eqref{eq:Jv4}, we require $\mu_{12}^{*}v_2 - \mu_{2}v_4 < 0$. Choose $v_4 = (1+\epsilon)\frac{\mu_{12}^{*}}{\mu_{2}}v_2$ for the same $\epsilon > 0$. Then,
\begin{equation*}
  (Jv)_4 = \mu_{12}^{*}v_2 - \mu_{2}(1+\epsilon)\frac{\mu_{12}^{*}}{\mu_{2}}v_2 = -\epsilon\mu_{12}^{*}v_2.
\end{equation*}
Since $v_2 > 0$ and $\epsilon > 0$, $(Jv)_4 < 0$.

Substitute these choices for $v_3$ and $v_4$ into \eqref{eq:Jv1} and \eqref{eq:Jv2}.
For \eqref{eq:Jv1}:
\begin{align*}
(Jv)_1 &= [(c_{12}-1)\theta_{11} - \mu_{11}^{*}]v_1 + c_{12}\theta_{12}v_2 + c_{12}\mu_{2}(1+\epsilon)\frac{\mu_{11}^{*}}{\mu_{2}}v_1 + c_{12}\mu_{2}(1+\epsilon)\frac{\mu_{12}^{*}}{\mu_{2}}v_2 \\
&= [(c_{12}-1)\theta_{11} - \mu_{11}^{*} + c_{12}(1+\epsilon)\mu_{11}^{*}]v_1 + [c_{12}\theta_{12} + c_{12}(1+\epsilon)\mu_{12}^{*}]v_2.
\end{align*}
Let $P_1 = \theta_{11} + \mu_{11}^{*} > 0$ and $P_2 = \theta_{12} + \mu_{12}^{*} > 0$. Then $(Jv)_1$ becomes:
\begin{equation}
(Jv)_1 = [(c_{12}-1)P_1 + c_{12}\epsilon\mu_{11}^{*}]v_1 + [c_{12}P_2 + c_{12}\epsilon\mu_{12}^{*}]v_2. \label{eq:Jv1_substituted}
\end{equation}
Similarly, $(Jv)_2$ becomes:
\begin{equation}
(Jv)_2 = [c_{22}P_1 + c_{22}\epsilon\mu_{11}^{*}]v_1 + [(c_{22}-1)P_2 + c_{22}\epsilon\mu_{12}^{*}]v_2. \label{eq:Jv2_substituted}
\end{equation}
We first choose $v_1, v_2 > 0$ such that the main terms (without $\epsilon$) in \eqref{eq:Jv1_substituted} and \eqref{eq:Jv2_substituted} are negative:
\begin{align}
(c_{12}-1)P_1 v_1 + c_{12}P_2 v_2 &< 0 \label{eq:cond_v1_main}, \\
c_{22}P_1 v_1 + (c_{22}-1)P_2 v_2 &< 0 \label{eq:cond_v2_main}.
\end{align}
Since $0 < c_{12} < 1$ and $0 < c_{22} < 1$ (as $q_i^*, z_0^* > 0$), we have $c_{12}-1 < 0$ and $c_{22}-1 < 0$.
From \eqref{eq:cond_v1_main}, $\frac{v_2}{v_1} < \frac{(1-c_{12})P_1}{c_{12}P_2}$.
From \eqref{eq:cond_v2_main}, $\frac{v_2}{v_1} > \frac{c_{22}P_1}{(1-c_{22})P_2}$.
A valid ratio $v_2/v_1$ exists if $\frac{c_{22}P_1}{(1-c_{22})P_2} < \frac{(1-c_{12})P_1}{c_{12}P_2}$, which simplifies to $\frac{c_{22}}{1-c_{22}} < \frac{1-c_{12}}{c_{12}}$ (since $P_1, P_2 > 0$).
This further simplifies to $c_{12}c_{22} < (1-c_{12})(1-c_{22})$, implying $c_{12} + c_{22} < 1$.
To show this:
\begin{equation*}
  c_{12} + c_{22} = \frac{q_1^*}{z_0^*+q_1^*+q_2^*} + \frac{q_2^*}{z_0^*+q_1^*+q_2^*} = \frac{q_1^*+q_2^*}{z_0^*+q_1^*+q_2^*}.
\end{equation*}
Since the equilibrium assumes active matching thresholds $\mu_{1i}^* > 0$, it implies $C (q_i^*)^{\alpha} (z_0^*)^{\alpha} = \mu_{1i}^* > 0$. This requires $q_i^* > 0$ and $z_0^* > 0$.
Therefore, the denominator $z_0^*+q_1^*+q_2^*$ is strictly greater than the numerator $q_1^*+q_2^*$ (as $z_0^* > 0$).
Thus, $\frac{q_1^*+q_2^*}{z_0^*+q_1^*+q_2^*} < 1$. This condition is satisfied.

Thus, we can choose $v_1, v_2 > 0$ such that for some $\delta_1, \delta_2 > 0$:
\begin{align*}
(c_{12}-1)P_1 v_1 + c_{12}P_2 v_2 &= -\delta_1, \\
c_{22}P_1 v_1 + (c_{22}-1)P_2 v_2 &= -\delta_2.
\end{align*}
Then, \eqref{eq:Jv1_substituted} and \eqref{eq:Jv2_substituted} become:
\begin{align*}
(Jv)_1 &= -\delta_1 + \epsilon [c_{12}\mu_{11}^{*}v_1 + c_{12}\mu_{12}^{*}v_2], \\
(Jv)_2 &= -\delta_2 + \epsilon [c_{22}\mu_{11}^{*}v_1 + c_{22}\mu_{12}^{*}v_2].
\end{align*}
Since $\delta_1, \delta_2 > 0$, and the terms multiplied by $\epsilon$ are positive (as $v_1,v_2>0$, $c_{ij} \ge 0$, $\mu_{ij}^* >0$), we can choose a sufficiently small $\epsilon > 0$ such that $(Jv)_1 < 0$ and $(Jv)_2 < 0$.
With such $v_1, v_2 > 0$ and $\epsilon > 0$, we have constructed $v \gg 0$ such that $Jv \ll 0$.
Therefore, all eigenvalues of $J$ have strictly negative real parts.
\textit{End of proof for eigenvalues of $J$.}

Since all eigenvalues of $J$ have strictly negative real parts, the solution $Z_{sub}(t)$ of the linear system $Z_{sub}' = J Z_{sub} + K$ converges exponentially fast to its unique equilibrium $Z_{sub}^* = -J^{-1}K = (z_{11}^*, z_{12}^*, z_{21}^*, z_{22}^*)$.

Combining parts (a) and (b): $z_0(t)$ and $q_i(t)$ converge to their equilibrium values. As they approach these values, the dynamics for $z_{1i}(t)$ and $z_{2i}(t)$ become those of a stable linear system, ensuring that these components also converge exponentially to their respective equilibrium values. This establishes the exponential convergence of the overall fluid process $(q_i(t), z_0(t), z_{1i}(t), z_{2i}(t))$ to the unique equilibrium point $X^*$.
\Halmos
\end{proof}

%\section*{EC.7. Justification of the (approximate) upper bound}
%
%We first restate the core optimization problem (27) as follows:
%\begin{equation}\label{eq:EC47}
%\max_{\mu_1} z_2,
%\end{equation}
%s.t.
%\begin{align}
%\lambda &= q\tau_0 + z_1\tau_1 + z_2\mu_2, \label{eq:EC48}\\
%z_1\mu_1 &= z_2\mu_2, \label{eq:EC49}\\
%1 &= z_1 + z_2 + z_0, \label{eq:EC50}\\
%\mu_1 &= C(q)^{\alpha_1} (z_0)^{\alpha_2}. \label{eq:EC51}
%\end{align}
%$q, z_0, z_1, z_2 \ge 0$.
%
%We denote the optimal value of Problem \eqref{eq:EC47} as $z_{22}(\lambda, \mu_2, C),$ reflecting the dependence on $\lambda, \mu_2, C.$ Next we establish a sensitivity result regarding $z_{22}(\lambda, \mu_2, C),$ which is central to our analysis later.
%
%\begin{lemma} \label{lem:EC3}
%Suppose Problem \eqref{eq:EC47} has feasible solutions under two sets of parameters, $(\lambda, \mu_2, C)$ and $(\lambda^*, \mu^*_2, C^*)$. Then
%\begin{equation}
%z_{22}(\lambda^*, \mu^*_2, C^*) \ge z_{22}(\lambda, \mu_2, C), \label{eq:EC52}
%\end{equation}
%when $(\lambda^*, \mu^*_2, C^*) \ge (\lambda, \mu_2, C).$
%\end{lemma}
%\begin{proof}
%For parameters $(\lambda, \mu_2, C)$, denote the optimal solution of Problem \eqref{eq:EC47} by $(\mu^2_1, q^2, z^2_0, z^2_1, z^2_2).$ We aim to find a feasible solution of Problem \eqref{eq:EC47} that is "close" to $(\mu^2_1, q^2, z^2_0, z^2_1, z^2_2)$ when the parameter set is changed to $(\lambda^*, \mu^*_2, C^*).$
%\end{proof}

